% =============================================================================
% Template Pandoc LaTeX pour la classe memoirerecherche
% Utilisé par bookdown::pdf_book via rapport.Rmd
% =============================================================================

% Pré-déclarer l'option [table] pour xcolor AVANT le chargement de la classe.
% memoirerecherche.cls fait \RequirePackage[table]{xcolor} ; si kableExtra
% (ou un autre paquet) tente ensuite \usepackage{xcolor}, LaTeX peut
% signaler un « Option clash for package xcolor ». La ligne suivante
% garantit la compatibilité dans tous les cas.
\PassOptionsToPackage{table}{xcolor}

\documentclass[11pt,a4paper]{memoirerecherche}

% --- Géométrie ----------------------------------------------------------------
\usepackage[left=2.5cm,right=2.5cm,top=2.5cm,bottom=2.5cm]{geometry}

% --- Encodage et polices (XeLaTeX) --------------------------------------------
\usepackage{fontspec}
\defaultfontfeatures{Ligatures=TeX,Scale=MatchLowercase}

% --- Langue -------------------------------------------------------------------
\usepackage[french]{babel}
\usepackage{csquotes}

% --- Mathématiques ------------------------------------------------------------
\usepackage{amsmath,amssymb}

% --- Tableaux et flottants ----------------------------------------------------
% Note : graphicx, xcolor, colortbl, hyperref, caption déjà chargés par memoirerecherche.cls
\usepackage{longtable}
\usepackage{booktabs}
\usepackage{array}
\usepackage{calc}   % Requis par Pandoc ≥ 3.x pour \real{} dans les largeurs de colonnes
\usepackage{multirow}
\usepackage{float}

% --- Graphiques (maxwidth/maxheight pour Pandoc) -------------------------------
\makeatletter
\def\maxwidth{\ifdim\Gin@nat@width>\linewidth\linewidth\else\Gin@nat@width\fi}
\def\maxheight{\ifdim\Gin@nat@height>\textheight\textheight\else\Gin@nat@height\fi}
\makeatother
\setkeys{Gin}{width=\maxwidth,height=\maxheight,keepaspectratio}

% --- Listes -------------------------------------------------------------------
\providecommand{\tightlist}{%
  \setlength{\itemsep}{0pt}\setlength{\parskip}{0pt}}

% --- Configuration hyperref (package déjà chargé par la classe) ---------------
\hypersetup{
  pdftitle={Analyse statistique de l'état de santé perçu en France :
modélisation multivariée, sélection de variables et
régression logistique ordinale à partir de l'Enquête
Santé et Protection Sociale (ESPS)},
  colorlinks=true,
  linkcolor=urcabrown,
  citecolor=urcabrown,
  urlcolor=urcalightbrown,
  breaklinks=true
}
\urlstyle{same}

% --- Numérotation des sections en « Chapitre N » -----------------------------
\renewcommand{\thesection}{\arabic{section}}
\renewcommand{\thesubsection}{\thesection.\arabic{subsection}}
\renewcommand{\thesubsubsection}{\thesubsection.\arabic{subsubsection}}

% Redéfinir le format des sections numérotées : « Chapitre N. Titre » en page nouvelle
\titleformat{\section}
  {\newpage\color{urcabrown}\normalfont\Large\bfseries}
  {\color{urcabrown}Chapitre~\thesection.}{1em}{}
\titlespacing*{\section}
  {0pt}{3.3ex}{3.3ex}

% Format des sections non-numérotées (Remerciements, Avant-propos, Conclusion, Annexes...)
\titleformat{name=\section,numberless}
  {\color{urcabrown}\normalfont\Large\bfseries}
  {}{0em}{}
\titlespacing*{name=\section,numberless}
  {0pt}{3.3ex}{3.3ex}

% --- En-têtes / pieds de page -------------------------------------------------
\usepackage{fancyhdr}
\pagestyle{fancy}
\fancyhf{}
\fancyhead[L]{\small\leftmark}
\fancyhead[R]{\small\thepage}
\renewcommand{\headrulewidth}{0.4pt}

% --- Bibliographie (biblatex) -------------------------------------------------
\usepackage[]{biblatex}
\addbibresource{biblio-sante.bib}

% --- Comptage automatique des références bibliographiques ---------------------
% Le compteur citnum (défini dans memoirerecherche.cls) est incrémenté pour
% chaque entrée imprimée dans la bibliographie ; \total{citnum} sur la page
% de titre affichera la valeur finale (disponible dès le second passage LaTeX).
\newcounter{citnum}
\AtEveryBibitem{\stepcounter{citnum}}

% --- Coloration syntaxique Pandoc ---------------------------------------------
% Pré-définir l'environnement Shaded (requis par les macros de coloration
% Pandoc qui utilisent \renewenvironment) et la couleur shadecolor.
\usepackage{tcolorbox}
\usepackage{framed}
\definecolor{shadecolor}{RGB}{248,248,248}
\newenvironment{Shaded}{\begin{snugshade}}{\end{snugshade}}

% --- Header-includes depuis le YAML -------------------------------------------
\setlength{\headheight}{24.80502pt}
\usepackage{underscore}
\usepackage{tcolorbox}
\tcbuselibrary{breakable}
\renewenvironment{Shaded}{\begin{tcolorbox}[colback=shadecolor, colframe=shadecolor, boxrule=0pt, sharp corners, breakable]}{\end{tcolorbox}}
\usepackage{booktabs}
\usepackage{longtable}
\usepackage{float}
\floatplacement{figure}{H}
\floatplacement{table}{H}
\usepackage{caption}
\captionsetup{font=small, labelfont=bf}
\usepackage{booktabs}
\usepackage{longtable}
\usepackage{array}
\usepackage{multirow}
\usepackage{wrapfig}
\usepackage{float}
\usepackage{colortbl}
\usepackage{pdflscape}
\usepackage{tabu}
\usepackage{threeparttable}
\usepackage{threeparttablex}
\usepackage[normalem]{ulem}
\usepackage{makecell}
\usepackage{xcolor}

% --- Variables personnalisées → commandes memoirerecherche.cls ----------------
\diplome{Master Statistique pour l'évaluation et la prévision, 1ère année}
\module{Modèles linéaires - Statistical Analysis System}
\enseig{CHARLOTTE PAPELARD}
\evaluation{Projet de Modèles linéaires avec Statistical Analysis System}
\coefficient{50\%}
\anac{2025-2026}

% Résumé
\resume{En France, 8,2~\% des adultes déclarent un état de santé « mauvais » ou « très mauvais », un chiffre qui atteint 19~\% parmi les ménages du premier quintile de revenus (DREES, 2017). Ce gradient social, documenté par la Commission OMS sur les déterminants sociaux de la santé (2008), soulève une question fondamentale : quels facteurs sociodémographiques, économiques et comportementaux déterminent la perception individuelle de l'état de santé, et comment les hiérarchiser de façon robuste ? Le présent projet met en œuvre une analyse statistique complète des déterminants de l'état de santé perçu en France métropolitaine, à partir de données issues de l'Enquête Santé et Protection Sociale (ESPS) de l'IRDES (500 individus, 25 variables). La variable réponse --- qualitative ordinale à cinq modalités (Très mauvais, Mauvais, Assez bon, Bon, Très bon) --- est analysée au moyen de sept familles de méthodes : analyse bivariée (\(\chi^2\), ANOVA, Kruskal-Wallis, Spearman), régression linéaire multiple sur le score ordinal (1--5), colinéarité (VIF, tolérance, indices de conditionnement), sélection de variables (\emph{forward}, \emph{backward}, \emph{stepwise}), validation complète (Shapiro-Wilk, Durbin-Watson, test de White, Cook, LOESS), ANCOVA (facteurs qualitatifs, moyennes marginales de Tukey), régression logistique ordinale (cotes proportionnelles, test du Score) et binaire (ROC, AUC, matrice de confusion). L'âge, les consultations médicales, les médicaments réguliers, l'activité physique et le handicap déclaré émergent comme déterminants robustes, tous significatifs au seuil de 5~\% à travers les différentes spécifications. Les limites --- hypothèse de proportionnalité des odds, absence de correction pour biais de sélection, taille de l'échantillon --- sont rigoureusement discutées.}

% Mots-clés
\keywords{État de santé perçu, Régression logistique ordinale, Régression linéaire multiple, ANCOVA, Sélection de variables, Enquête ESPS, Colinéarité, Déterminants sociaux de la santé}

% Dépôt Git et base de données
\gitrepo{\url{https://gitlab-mi.univ-reims.fr/adja0021}}
\database{ESPS : base\_esps.csv (500 obs., 25 variables) --- IRDES}

% Logo de la page de titre
\logo{logoSESGURCA.jpg}

% Métriques de la page de titre — comptage automatique
% Les compteurs figure, table et citnum sont gérés automatiquement par
% totcount (voir memoirerecherche.cls). Les lignes ci-dessous ne sont
% conservées que pour ne pas casser un éventuel code qui référencerait
% \@totfig, \@tottab, etc.
\totfig{\totalfigures}\tottab{\totaltables}\totref{--}\tottheo{0}

% --- Titre, auteur, date ------------------------------------------------------
\title{Analyse statistique de l'état de santé perçu en France :
modélisation multivariée, sélection de variables et
régression logistique ordinale à partir de l'Enquête
Santé et Protection Sociale (ESPS)}
% N.B. : ne PAS utiliser \and ici — la classe memoirerecherche.cls affiche
% \@author dans un minipage/tabular ; or, \and est défini par article.cls
% comme \end{tabular}\begin{tabular}, ce qui casse l'imbrication de tables
% du \maketitle de la classe et provoque « ! Missing } inserted. ».
\author{Hervé ADJAKOSSA}
\date{04 May 2026}

% =============================================================================
\begin{document}

\maketitle
\newpage

% --- Table des matières -------------------------------------------------------
{
\setcounter{tocdepth}{3}
\tableofcontents
\newpage
}

% --- Liste des tables ---------------------------------------------------------

% --- Liste des figures --------------------------------------------------------

% --- Corps du document --------------------------------------------------------
\newpage

\thispagestyle{empty}
\vspace*{\fill}
\begin{flushright}
\textit{«\ La santé est un état de complet bien-être physique,\\
mental et social, et ne consiste pas seulement\\
en une absence de maladie ou d'infirmité.\ »}

\medskip
--- Organisation mondiale de la santé, 1946
\end{flushright}
\vspace{1cm}
\begin{center}
\textit{Je dédie ce travail à ma famille et à tous ceux qui\\
œuvrent pour la réduction des inégalités de santé.}
\end{center}
\vspace*{\fill}

\newpage

\hypertarget{remerciements}{%
\section*{Remerciements}\label{remerciements}}
\addcontentsline{toc}{section}{Remerciements}

Je tiens à exprimer ma profonde gratitude à \textbf{Charlotte Papelard}, responsable du module \emph{Modèles linéaires --- Statistical Analysis System}, pour la qualité de son enseignement, ses conseils méthodologiques et l'exigence scientifique qu'elle a su transmettre tout au long de cette formation. Son accompagnement a été déterminant dans la conduite de ce projet.

Mes remerciements vont également à l'ensemble des enseignants du Master \emph{Statistique pour l'évaluation et la prévision} de l'Université de Reims Champagne-Ardenne, dont les enseignements ont fourni les fondements théoriques et pratiques mobilisés dans ce travail.

Je remercie l'Institut de recherche et documentation en économie de la santé (IRDES) pour la mise à disposition des données de l'Enquête Santé et Protection Sociale (ESPS), sans lesquelles cette analyse n'aurait pu être réalisée.

Enfin, je tiens à remercier ma famille et mes proches pour leur soutien constant et leurs encouragements tout au long de cette année universitaire.

\newpage

\hypertarget{avant-propos}{%
\section*{Avant-propos}\label{avant-propos}}
\addcontentsline{toc}{section}{Avant-propos}

\hypertarget{contexte-et-objectifs}{%
\subsection*{Contexte et objectifs}\label{contexte-et-objectifs}}
\addcontentsline{toc}{subsection}{Contexte et objectifs}

En France, les inégalités sociales de santé constituent un enjeu majeur de santé publique. La Direction de la recherche, des études, de l'évaluation et des statistiques (DREES) rapporte que 8,2~\% des adultes déclarent un état de santé « mauvais » ou « très mauvais », un chiffre qui atteint 19~\% parmi les ménages du premier quintile de revenus. Ce gradient social, documenté par la Commission de l'OMS sur les déterminants sociaux de la santé (2008), soulève une question fondamentale : quels facteurs sociodémographiques, économiques et comportementaux déterminent la perception individuelle de l'état de santé, et comment les hiérarchiser de façon robuste ?

Le présent projet met en œuvre une analyse statistique complète des déterminants de l'état de santé perçu en France métropolitaine, à partir de données issues de l'Enquête Santé et Protection Sociale (ESPS) de l'IRDES (500 individus, 25 variables). La variable réponse --- qualitative ordinale à cinq modalités (Très mauvais, Mauvais, Assez bon, Bon, Très bon) --- est analysée au moyen de sept familles de méthodes : analyse bivariée, régression linéaire multiple, colinéarité, sélection de variables, validation complète, ANCOVA et régression logistique.

\hypertarget{plan-du-rapport}{%
\subsection*{Plan du rapport}\label{plan-du-rapport}}
\addcontentsline{toc}{subsection}{Plan du rapport}

Ce rapport s'organise autour de sept chapitres. Chaque chapitre s'ouvre sur un énoncé introductif des différents points qui y sont abordés et se conclut par une synthèse.

Le \textbf{chapitre 1} est une \textbf{introduction générale} qui présente le contexte, la problématique, les hypothèses de recherche et la stratégie méthodologique retenue.

Le \textbf{chapitre 2} pose le \textbf{cadre conceptuel, théorique et analytique} en justifiant chaque choix méthodologique au regard de la littérature sur la santé perçue et les déterminants sociaux.

Le \textbf{chapitre 3} propose une \textbf{revue de littérature} sur les déterminants de la santé perçue et les méthodes de régression appliquées à ce domaine, en identifiant la contribution spécifique de notre travail.

Le \textbf{chapitre 4} détaille la \textbf{méthodologie} statistique retenue : design de recherche, spécification des modèles (régression linéaire, ANCOVA, logistique ordinale et binaire) et stratégie de validation.

Le \textbf{chapitre 5} présente les \textbf{données et la construction de la base empirique} : sources, construction des variables et statistiques descriptives commentées.

Le \textbf{chapitre 6} expose les \textbf{résultats} : analyses bivariées et multivariées, diagnostics de validation, régression logistique ordinale et binaire.

Le \textbf{chapitre 7} discute la \textbf{robustesse, les limites et les perspectives} du travail, en confrontant systématiquement les résultats à la littérature épidémiologique.

Le rapport se termine par une \textbf{conclusion générale} structurée en bilan et perspectives, suivie d'\textbf{annexes} comprenant le dictionnaire de données et le code SAS complet.

\newpage

\hypertarget{sigles-et-acronymes}{%
\section*{Sigles et acronymes}\label{sigles-et-acronymes}}
\addcontentsline{toc}{section}{Sigles et acronymes}

\begin{tabular}{@{} l @{\hspace{1.5em}} p{12cm} @{}}
\toprule
\textbf{Sigle} & \textbf{Signification} \\
\midrule
AIC & \textit{Akaike Information Criterion} --- Critère d'information d'Akaike \\
ANCOVA & \textit{Analysis of Covariance} --- Analyse de covariance \\
ANOVA & \textit{Analysis of Variance} --- Analyse de la variance \\
AUC & \textit{Area Under the Curve} --- Aire sous la courbe ROC \\
BIC & \textit{Bayesian Information Criterion} --- Critère d'information bayésien \\
CMU-C & Couverture maladie universelle complémentaire \\
DREES & Direction de la recherche, des études, de l'évaluation et des statistiques \\
ESPS & Enquête Santé et Protection Sociale \\
EU-SILC & \textit{European Union --- Statistics on Income and Living Conditions} \\
GLM & \textit{General Linear Model} --- Modèle linéaire général \\
IMC & Indice de masse corporelle \\
INSEE & Institut national de la statistique et des études économiques \\
IRDES & Institut de recherche et documentation en économie de la santé \\
LOESS & \textit{Locally Estimated Scatterplot Smoothing} --- Lissage local \\
LSMEANS & \textit{Least Squares Means} --- Moyennes des moindres carrés \\
MCO & Moindres carrés ordinaires \\
OMS & Organisation mondiale de la santé \\
ROC & \textit{Receiver Operating Characteristic} --- Courbe caractéristique de fonctionnement du récepteur \\
SAS & \textit{Statistical Analysis System} \\
SHARE & \textit{Survey of Health, Ageing and Retirement in Europe} \\
SRH & \textit{Self-Rated Health} --- Santé auto-évaluée \\
VIF & \textit{Variance Inflation Factor} --- Facteur d'inflation de la variance \\
\bottomrule
\end{tabular}

\newpage

\hypertarget{chap-introduction}{%
\section{Introduction générale}\label{chap-introduction}}

\hypertarget{accroche-et-contexte-guxe9nuxe9ral}{%
\subsection{Accroche et contexte général}\label{accroche-et-contexte-guxe9nuxe9ral}}

En 2019, 67,4~\% des Français déclaraient un état de santé « bon » ou « très bon », un chiffre qui masque des disparités considérables : parmi les ménages du premier quintile de revenus, la proportion d'adultes se déclarant en « mauvaise » ou « très mauvaise » santé atteint 19~\%, soit près du triple de la moyenne nationale\footnote{Direction de la recherche, des études, de l'évaluation et des statistiques (DREES), \emph{L'état de santé de la population en France --- Rapport 2017}, Paris.}. Ce constat, loin d'être anecdotique, traduit un phénomène structurel : les inégalités sociales de santé perçue s'inscrivent au cœur du débat économique et sanitaire contemporain. L'Organisation mondiale de la santé (OMS), à travers sa Commission sur les déterminants sociaux de la santé, a établi que « les conditions dans lesquelles les individus naissent, grandissent, vivent, travaillent et vieillissent constituent les déterminants sociaux de la santé, lesquels sont principalement responsables des inégalités en santé »\footnote{\textcite{oms2008}, Commission des déterminants sociaux de la santé, Rapport final, Genève.}. Ce constat appelle une quantification rigoureuse des facteurs qui façonnent la perception individuelle de l'état de santé.

La France dispose, à cet égard, d'un instrument d'observation privilégié : l'Enquête Santé et Protection Sociale (ESPS), conduite depuis 1988 par l'Institut de recherche et documentation en économie de la santé (IRDES)\footnote{\textcite{irdes_esps}.}. Cette enquête, qui articule données de santé déclarée, recours aux soins et caractéristiques socio-économiques, offre un terrain empirique particulièrement riche pour l'analyse des déterminants de l'état de santé perçu --- variable centrale de notre étude.

\hypertarget{intuxe9ruxeat-scientifique-et-empirique-pourquoi-uxe9tudier-la-santuxe9-peruxe7ue-par-la-ruxe9gression-linuxe9aire}{%
\subsection{Intérêt scientifique et empirique : pourquoi étudier la santé perçue par la régression linéaire ?}\label{intuxe9ruxeat-scientifique-et-empirique-pourquoi-uxe9tudier-la-santuxe9-peruxe7ue-par-la-ruxe9gression-linuxe9aire}}

L'état de santé perçu (\emph{self-rated health}, SRH) occupe une place singulière en épidémiologie et en économie de la santé. \textcite{idler1997}, dans une méta-analyse pionnière portant sur vingt-sept études de cohorte, ont démontré que cette mesure subjective constitue un prédicteur indépendant de la mortalité, y compris après ajustement sur l'état de santé objectif. \textcite{jylha2009} propose un modèle conceptuel unifié expliquant ce pouvoir prédictif : la santé perçue intégrerait simultanément des informations biologiques infracliniques, des comparaisons sociales et des représentations culturelles de la santé\footnote{Ce résultat explique pourquoi les grandes enquêtes de population --- EU-SILC, ESPS, SHARE --- incluent systématiquement une question d'auto-évaluation de la santé.}.

La régression linéaire multiple, bien que la variable réponse soit ordinale, est largement mobilisée dans la littérature pour modéliser la santé perçue. \textcite{kennedy1998} utilisent la régression linéaire sur un score de santé auto-évalué pour démontrer l'effet des inégalités de revenus sur la santé perçue aux États-Unis. \textcite{subramanian2002} recourent à des modèles linéaires multiniveaux pour estimer l'effet contextuel du capital social sur la santé auto-déclarée. En France, \textcite{devaux2010} et \textcite{etile2007} mobilisent la régression linéaire pour identifier l'hétérogénéité de déclaration selon le statut socio-économique. Cette approche, qui suppose l'équidistance entre les modalités ordinales, présente l'avantage de la parcimonie et de l'interprétabilité directe des coefficients --- un atout essentiel pour l'aide à la décision en santé publique\footnote{La régression logistique ordinale, plus rigoureuse sur le plan théorique, sera néanmoins estimée à titre de comparaison et de robustesse.}.

\textbf{Notre motivation} pour ce choix méthodologique est triple : (i) reproduire et discuter les résultats de la littérature dans le cadre pédagogique de l'ESPS ; (ii) confronter systématiquement la régression linéaire à la régression logistique ordinale pour évaluer la robustesse des conclusions ; (iii) mobiliser un arsenal complet de diagnostics statistiques (colinéarité, validation, sélection) rarement déployés conjointement dans les études existantes.

\hypertarget{probluxe9matique-guxe9nuxe9rale}{%
\subsection{Problématique générale}\label{probluxe9matique-guxe9nuxe9rale}}

L'enjeu de la modélisation de la santé perçue se situe à l'intersection de deux tensions. D'une part, la nature ordinale de la variable réponse interroge la légitimité de la régression linéaire : l'hypothèse d'équidistance entre « Très mauvais » et « Mauvais » est-elle soutenable ? D'autre part, la multiplicité des déterminants potentiels --- biologiques, socio-économiques, comportementaux, institutionnels --- engendre des effets de confusion qui requièrent une stratégie d'identification rigoureuse\footnote{L'endogénéité potentielle des variables de recours aux soins (consultations, médicaments) constitue une difficulté supplémentaire, car ces variables sont simultanément cause et conséquence de l'état de santé (\textcite{grossman1972}).}.

Comment, dès lors, modéliser et hiérarchiser les déterminants de l'état de santé perçu en France, en tenant compte de la nature ordinale de la variable réponse, des effets de confusion entre variables sociodémographiques, économiques et comportementales, et des limites inhérentes aux données d'enquête en coupe transversale ?

\hypertarget{question-centrale-de-recherche}{%
\subsection{Question centrale de recherche}\label{question-centrale-de-recherche}}

\textbf{Quels sont les facteurs sociodémographiques, économiques et comportementaux qui déterminent significativement l'état de santé perçu des individus en France métropolitaine, et dans quelle mesure la régression logistique ordinale offre-t-elle une modélisation plus adéquate que la régression linéaire sur le score ordinal ?}

\hypertarget{objectif-guxe9nuxe9ral-et-objectifs-spuxe9cifiques}{%
\subsection{Objectif général et objectifs spécifiques}\label{objectif-guxe9nuxe9ral-et-objectifs-spuxe9cifiques}}

\textbf{Objectif général} : Identifier, hiérarchiser et quantifier les déterminants de l'état de santé perçu en France à partir des données de l'ESPS, en mobilisant un arsenal complet de méthodes statistiques et en confrontant systématiquement les résultats à la réalité épidémiologique documentée.

\textbf{Objectifs spécifiques} :

\begin{enumerate}
\def\labelenumi{\arabic{enumi}.}
\tightlist
\item
  \textbf{Caractériser les associations bivariées} entre l'état de santé perçu et les variables explicatives (qualitatives et quantitatives), en justifiant le choix de chaque test statistique par la nature des variables croisées.
\item
  \textbf{Construire et valider un modèle de régression linéaire multiple} sur le score ordinal via \texttt{PROC\ REG}, en étudiant la colinéarité (VIF, indices de conditionnement) et en sélectionnant les variables pertinentes par critères informationnels.
\item
  \textbf{Estimer un modèle ANCOVA} via \texttt{PROC\ GLM} combinant covariables quantitatives et facteurs qualitatifs, avec tests d'interaction et moyennes marginales (Tukey/LSMEANS), pour isoler l'effet propre des déterminants comportementaux.
\item
  \textbf{Estimer un modèle de régression logistique ordinale} via \texttt{PROC\ LOGISTIC} (cotes proportionnelles) et un modèle logistique binaire, évaluer leur performance prédictive (ROC, AUC, matrice de confusion) et comparer les approches.
\item
  \textbf{Évaluer la capacité prédictive} des modèles retenus et confronter les résultats aux connaissances empiriques existantes pour juger de la conformité à la réalité.
\end{enumerate}

\hypertarget{hypothuxe8ses-de-recherche}{%
\subsection{Hypothèses de recherche}\label{hypothuxe8ses-de-recherche}}

\begin{itemize}
\tightlist
\item
  \textbf{H1} : L'âge, le nombre de consultations médicales et le nombre de médicaments réguliers sont les déterminants quantitatifs les plus fortement associés à l'état de santé perçu --- conformément au modèle de capital santé de \textcite{grossman1972}, selon lequel l'état de santé se déprécie avec l'âge et la consommation médicale reflète le stock de santé résiduel.
\item
  \textbf{H2} : L'activité physique et le handicap déclaré exercent un effet significatif sur l'état de santé perçu, au-delà des variables sociodémographiques et économiques (effet propre mis en évidence par l'ANCOVA) --- ce qui prolonge les résultats de \textcite{singh2006} sur les déterminants comportementaux de la santé perçue.
\item
  \textbf{H3} : Le modèle logistique ordinal, respectant la nature qualitative ordinale de la variable réponse, offre un ajustement supérieur au modèle linéaire sur le score numérique --- hypothèse fondée sur les travaux de \textcite{mccullagh1980} et \textcite{agresti2010} qui établissent la supériorité théorique du \emph{cumulative logit model} pour les réponses ordinales.
\end{itemize}

\hypertarget{positionnement-scientifique}{%
\subsection{Positionnement scientifique}\label{positionnement-scientifique}}

Ce travail s'inscrit au carrefour de la biostatistique, de l'économie de la santé et de l'économétrie appliquée. Il mobilise des méthodes de régression en grande dimension (sélection de variables, diagnostics de colinéarité) et de régression logistique ordinale, appliquées aux données de l'ESPS. Le positionnement est à la fois \textbf{empirique} --- identifier les déterminants de la santé perçue en France --- et \textbf{méthodologique} --- confronter les approches linéaire et ordinale pour évaluer la robustesse des conclusions.

\hypertarget{stratuxe9gie-de-recherche-et-duxe9marche-muxe9thodologique}{%
\subsection{Stratégie de recherche et démarche méthodologique}\label{stratuxe9gie-de-recherche-et-duxe9marche-muxe9thodologique}}

La démarche suivie est séquentielle et cumulative, articulée en trois phases :

\begin{enumerate}
\def\labelenumi{\arabic{enumi}.}
\tightlist
\item
  \textbf{Phase exploratoire} : analyse bivariée (\texttt{PROC\ FREQ} pour les tests du \(\chi^2\), \texttt{PROC\ GLM} pour l'ANOVA, \texttt{PROC\ NPAR1WAY} pour le test de Kruskal-Wallis, \texttt{PROC\ CORR} pour la corrélation de Spearman) pour identifier les associations brutes entre l'état de santé perçu et chaque variable explicative, sans contrôle des effets confondants.
\item
  \textbf{Phase de modélisation} : régression linéaire multiple via \texttt{PROC\ REG} avec diagnostics complets (colinéarité via options VIF/TOL/COLLIN, validation des hypothèses via SPEC/DWPROB), sélection de variables par critères informationnels (sélection ADJRSQ, Forward, Backward, Stepwise), estimation ANCOVA via \texttt{PROC\ GLM} avec LSMEANS.
\item
  \textbf{Phase de confrontation et robustesse} : régression logistique ordinale et binaire via \texttt{PROC\ LOGISTIC} avec évaluation prédictive (courbe ROC, AUC, table de classement CTABLE), test du Score pour la proportionnalité des cotes, et discussion critique des résultats au regard de la littérature.
\end{enumerate}

Le terrain d'application est la France métropolitaine, cadre géographique justifié par la disponibilité de l'ESPS et par la richesse de la littérature comparative (gradient social documenté par \textcite{van_doorslaer2004} ; \textcite{art2006}).

\hypertarget{annonce-du-plan}{%
\subsection{Annonce du plan}\label{annonce-du-plan}}

Le \textbf{chapitre~\ref{chap-cadre}} pose le cadre conceptuel, théorique et analytique en justifiant chaque choix méthodologique. Le \textbf{chapitre~\ref{chap-litterature}} propose un état de l'art sur les déterminants de la santé perçue et les méthodes de régression linéaire appliquées à ce domaine, en identifiant la contribution spécifique de notre travail. Le \textbf{chapitre~\ref{chap-methodologie}} détaille la méthodologie statistique retenue et justifie la légitimité de chaque test. Le \textbf{chapitre~\ref{chap-donnees}} présente les données, la construction des variables et les statistiques descriptives commentées. Le \textbf{chapitre~\ref{chap-resultats}} expose l'ensemble des résultats, chaque analyse étant interprétée au regard de ses implications sanitaires et économiques. Le \textbf{chapitre~\ref{chap-robustesse}} discute la robustesse, les limites et les perspectives.

\hypertarget{chap-cadre}{%
\section{Cadre conceptuel, théorique et analytique}\label{chap-cadre}}

Ce chapitre pose les fondements conceptuels et théoriques de notre démarche. Chaque choix --- du cadre d'analyse au type de modèle --- est explicitement justifié au regard de la littérature et de la nature des données mobilisées.

\hypertarget{cadre-conceptuel-santuxe9-peruxe7ue-et-duxe9terminants-sociaux}{%
\subsection{Cadre conceptuel : santé perçue et déterminants sociaux}\label{cadre-conceptuel-santuxe9-peruxe7ue-et-duxe9terminants-sociaux}}

L'état de santé perçu (\emph{self-rated health}, SRH) est une mesure subjective globale où l'individu évalue sa propre santé sur une échelle ordinale (Très bon \(\rightarrow\) Très mauvais). Cette mesure, malgré sa simplicité apparente, capture des dimensions multiples de la santé --- physique, mentale, fonctionnelle --- et s'est révélée un prédicteur indépendant de la mortalité dans de nombreuses cohortes prospectives \autocite{idler1997}\footnote{Le pouvoir prédictif de la santé perçue a été confirmé dans plus de 27 études de cohorte, couvrant des contextes culturels variés (États-Unis, Europe, Japon).}.

\textbf{Pourquoi retenir la santé perçue comme variable réponse ?} Trois arguments justifient ce choix : (i) la santé perçue est disponible dans toutes les grandes enquêtes de population (ESPS, EU-SILC, SHARE), ce qui autorise les comparaisons internationales \autocite{van_doorslaer2004} ; (ii) elle agrège des informations biologiques et psychosociales inaccessibles par les seuls indicateurs objectifs \autocite{jylha2009} ; (iii) elle est corrélée aux comportements de santé et au recours aux soins, ce qui en fait un outil pertinent pour l'évaluation des politiques de santé publique\footnote{\textcite{singh2006} montrent, sur les cohortes Whitehall II et Gazel, que la santé perçue reflète des processus physiologiques infracliniques détectables par des marqueurs biologiques.}.

Les déterminants sociaux de la santé, selon le modèle de \textcite{dahlgren1991}, s'organisent en couches concentriques : facteurs biologiques individuels (âge, sexe), comportements et modes de vie (tabagisme, activité physique, alimentation), réseaux sociaux, conditions de vie et de travail, et enfin conditions socio-économiques, culturelles et environnementales générales. Ce modèle constitue notre grille de lecture pour la sélection des variables explicatives.

\hypertarget{cadre-thuxe9orique-moduxe8les-explicatifs-de-la-santuxe9-peruxe7ue}{%
\subsection{Cadre théorique : modèles explicatifs de la santé perçue}\label{cadre-thuxe9orique-moduxe8les-explicatifs-de-la-santuxe9-peruxe7ue}}

Plusieurs cadres théoriques sous-tendent l'analyse des déterminants de la santé perçue :

\begin{itemize}
\tightlist
\item
  \textbf{Modèle de capital santé de Grossman} \autocite{grossman1972} : l'état de santé est un stock qui se déprécie avec l'âge et dans lequel l'individu investit (consultations, médicaments, activité physique). Ce modèle justifie l'inclusion de l'âge, des consultations et des médicaments comme déterminants centraux\footnote{Le modèle de Grossman prédit que la demande de santé croît avec le niveau d'éducation, car l'éducation augmente l'efficience de la production de santé.}.
\item
  \textbf{Modèle socio-économique} \autocite{marmot2005} : les inégalités de revenus, d'éducation et de profession génèrent des gradients de santé perçue. Ce cadre justifie l'inclusion des variables de revenus, de profession et de couverture assurantielle.
\item
  \textbf{Modèle comportemental} : les comportements à risque (tabagisme, sédentarité, consommation d'alcool) médiatisent partiellement les effets socio-économiques sur la santé \autocite{bricard2015}. Ceci motive l'inclusion des variables de tabagisme, d'activité physique et de consommation d'alcool.
\item
  \textbf{Modèle d'accès aux soins} \autocite{despres2011} : la couverture assurantielle et le renoncement aux soins modulent la relation entre besoin de santé et état perçu, justifiant l'inclusion de la couverture complémentaire et du renoncement aux soins.
\end{itemize}

\hypertarget{cadre-analytique-justification-des-muxe9thodes-statistiques-retenues}{%
\subsection{Cadre analytique : justification des méthodes statistiques retenues}\label{cadre-analytique-justification-des-muxe9thodes-statistiques-retenues}}

\hypertarget{traitement-de-la-variable-ruxe9ponse-ordinale}{%
\subsubsection{Traitement de la variable réponse ordinale}\label{traitement-de-la-variable-ruxe9ponse-ordinale}}

La variable \texttt{etat\_sante\_percu} (5 modalités ordonnées) pose un choix méthodologique fondamental, que nous tranchons par une approche duale :

\begin{itemize}
\tightlist
\item
  \textbf{Approche numérique} : conversion en score (1--5) et utilisation des méthodes de régression linéaire classique (MCO via \texttt{PROC\ REG}, ANCOVA via \texttt{PROC\ GLM}). Cette approximation, bien que reposant sur l'hypothèse d'équidistance entre les modalités, est défendue par \textcite{wooldridge2015} dans les cas où l'échelle ordinale comporte au moins cinq modalités.
\item
  \textbf{Approche ordinale} : régression logistique ordinale (\emph{cumulative logit model}, modèle à cotes proportionnelles) via \texttt{PROC\ LOGISTIC} avec l'option \texttt{link=cumlogit}, qui respecte la nature qualitative de la réponse.
\end{itemize}

\textbf{Pourquoi combiner les deux approches ?} La confrontation systématique permet de (i) évaluer la robustesse des résultats à la spécification du modèle ; (ii) identifier les variables dont l'effet est sensible au traitement de l'échelle ordinale ; (iii) apporter une contribution méthodologique rarement entreprise sur un même jeu de données.

\hypertarget{limites-structurelles-du-design-choisi}{%
\subsubsection{Limites structurelles du design choisi}\label{limites-structurelles-du-design-choisi}}

\begin{itemize}
\tightlist
\item
  \textbf{Causalité} : l'enquête en coupe transversale ne permet pas d'établir de lien causal entre les déterminants et l'état de santé perçu. Les associations identifiées sont conditionnelles, non causales.
\item
  \textbf{Endogénéité} : certaines variables explicatives (consultations, médicaments) sont potentiellement endogènes, car déterminées simultanément par l'état de santé \autocite{grossman1972}.
\end{itemize}

\bigskip

\textbf{Conclusion du chapitre.} Le cadre analytique articule trois niveaux de justification : conceptuel (modèle de Dahlgren et Whitehead), théorique (capital santé de Grossman, gradient social de Marmot) et méthodologique (approche duale linéaire/ordinale, diagnostics complets). Les traitements statistiques sont réalisés à l'aide du logiciel SAS.

\hypertarget{chap-litterature}{%
\section{Revue de littérature}\label{chap-litterature}}

Ce chapitre dresse un état de l'art sur les déterminants de la santé perçue et les méthodes de régression appliquées à ce domaine.

\hypertarget{mesures-subjectives-de-la-santuxe9-fondements-et-validituxe9}{%
\subsection{Mesures subjectives de la santé : fondements et validité}\label{mesures-subjectives-de-la-santuxe9-fondements-et-validituxe9}}

L'état de santé perçu (\emph{self-rated health}) est la mesure subjective de santé la plus utilisée en épidémiologie sociale et en économie de la santé. \textcite{idler1997} ont établi, dans une méta-analyse portant sur vingt-sept études de cohorte, que l'auto-évaluation de la santé est un prédicteur indépendant de la mortalité --- résultat confirmé par \textcite{desalvo2006} dans une synthèse portant sur plus de 22 études longitudinales. \textcite{jylha2009} propose un modèle intégratif dans lequel la santé perçue résulte de trois processus cognitifs : l'évaluation du fonctionnement corporel, la comparaison sociale et la représentation culturelle de la santé.

\hypertarget{duxe9terminants-socio-uxe9conomiques-de-la-santuxe9-peruxe7ue}{%
\subsection{Déterminants socio-économiques de la santé perçue}\label{duxe9terminants-socio-uxe9conomiques-de-la-santuxe9-peruxe7ue}}

\hypertarget{effets-de-luxe2ge-et-du-sexe}{%
\subsubsection{Effets de l'âge et du sexe}\label{effets-de-luxe2ge-et-du-sexe}}

\textcite{mackenbach1999} documentent un gradient âge--santé perçue quasi universel : la proportion d'individus se déclarant en mauvaise santé croît linéairement avec l'âge dans l'ensemble des pays européens étudiés. \textcite{bambra2009} mettent en évidence un écart de genre persistant : les femmes déclarent systématiquement une moins bonne santé que les hommes, même après ajustement sur les conditions objectives de santé.

\hypertarget{gradient-social-de-santuxe9}{%
\subsubsection{Gradient social de santé}\label{gradient-social-de-santuxe9}}

Le gradient social de santé perçue est l'un des résultats les plus robustes de l'épidémiologie sociale. \textcite{marmot2005}, à partir des études Whitehall, montre que chaque échelon de la hiérarchie socio-professionnelle est associé à un état de santé perçu statistiquement inférieur à celui de l'échelon supérieur. \textcite{van_doorslaer2004}, dans une comparaison internationale portant sur treize pays européens, identifient la France comme l'un des pays présentant les inégalités sociales de santé perçue les plus marquées.

\hypertarget{ruxf4le-de-la-couverture-assurantielle-et-du-renoncement-aux-soins}{%
\subsubsection{Rôle de la couverture assurantielle et du renoncement aux soins}\label{ruxf4le-de-la-couverture-assurantielle-et-du-renoncement-aux-soins}}

\textcite{despres2011} montrent que le renoncement aux soins --- phénomène touchant 15 à 25~\% de la population selon les enquêtes --- est associé à un état de santé perçu significativement dégradé. La couverture complémentaire joue un rôle protecteur, en réduisant le reste à charge et en facilitant l'accès aux soins de premier recours \autocite{dourgnon2012}.

\hypertarget{la-ruxe9gression-linuxe9aire-pour-expliquer-la-santuxe9-peruxe7ue}{%
\subsection{La régression linéaire pour expliquer la santé perçue}\label{la-ruxe9gression-linuxe9aire-pour-expliquer-la-santuxe9-peruxe7ue}}

La conversion de la santé perçue en score numérique (1--5) pour application de la régression linéaire est une pratique courante en économie de la santé. \textcite{kennedy1998} ont ouvert la voie en utilisant la régression linéaire sur un score de santé auto-évalué. En France, \textcite{devaux2010} et \textcite{etile2007} mobilisent la régression linéaire pour identifier l'hétérogénéité de déclaration selon le statut socio-économique. L'hypothèse d'équidistance est contestée par \textcite{boes2006}, mais \textcite{greene2018} souligne que pour des échelles à 5 modalités, la régression linéaire fournit des approximations raisonnables.

\hypertarget{positionnement-de-la-contribution}{%
\subsection{Positionnement de la contribution}\label{positionnement-de-la-contribution}}

Ce travail se distingue par la mise en œuvre systématique et comparée de l'ensemble des méthodes (régression linéaire, ANCOVA, régression logistique ordinale et binaire) sur un même jeu de données, en utilisant SAS comme outil d'analyse.

\bigskip

\textbf{Conclusion du chapitre.} L'état de l'art révèle que la régression linéaire constitue un outil pertinent --- mais non exclusif --- pour modéliser la santé perçue. Notre contribution réside dans la confrontation systématique des approches linéaire et ordinale sur un même terrain empirique.

\hypertarget{chap-methodologie}{%
\section{Méthodologie}\label{chap-methodologie}}

Ce chapitre détaille la méthodologie statistique retenue. Les traitements statistiques sont réalisés à l'aide du logiciel SAS (Statistical Analysis System), outil de référence en biostatistique et en économétrie.

\hypertarget{design-de-recherche-terrain-et-pertinence}{%
\subsection{Design de recherche : terrain et pertinence}\label{design-de-recherche-terrain-et-pertinence}}

L'analyse repose sur des données en coupe transversale issues de l'Enquête Santé et Protection Sociale, conduite en \textbf{France métropolitaine}. Ce périmètre géographique est justifié par l'existence de l'ESPS, la richesse de la littérature comparative, et la spécificité du système de santé français (couverture universelle avec reste à charge modulé par la complémentaire).

\hypertarget{spuxe9cification-des-moduxe8les}{%
\subsection{Spécification des modèles}\label{spuxe9cification-des-moduxe8les}}

\hypertarget{moduxe8le-de-ruxe9gression-linuxe9aire-multiple}{%
\subsubsection{Modèle de régression linéaire multiple}\label{moduxe8le-de-ruxe9gression-linuxe9aire-multiple}}

Le score ordinal \texttt{score\_sante} (\(Y \in \{1, 2, 3, 4, 5\}\)) est modélisé comme :

\[Y_i = \beta_0 + \beta_1 \cdot \text{age}_i + \beta_2 \cdot \text{consultations}_i + \beta_3 \cdot \text{medicaments}_i + \cdots + \varepsilon_i\]

avec \(\varepsilon_i \sim \mathcal{N}(0, \sigma^2)\). Le modèle est estimé par \texttt{PROC\ REG} sous SAS.

\hypertarget{moduxe8le-ancova}{%
\subsubsection{Modèle ANCOVA}\label{moduxe8le-ancova}}

Le modèle ANCOVA étend la régression linéaire en incluant des facteurs qualitatifs via \texttt{PROC\ GLM} avec l'option \texttt{/SS3} pour la table ANOVA de type III.

\hypertarget{moduxe8le-logistique-ordinal-cotes-proportionnelles}{%
\subsubsection{Modèle logistique ordinal (cotes proportionnelles)}\label{moduxe8le-logistique-ordinal-cotes-proportionnelles}}

\[\log\left(\frac{P(Y \leq j)}{P(Y > j)}\right) = \alpha_j - \mathbf{X}' \boldsymbol{\beta}, \quad j = 1, \ldots, J-1\]

Estimé via \texttt{PROC\ LOGISTIC} avec l'option \texttt{link=cumlogit}. Le test du Score vérifie l'hypothèse de proportionnalité.

\hypertarget{moduxe8le-logistique-binaire}{%
\subsubsection{Modèle logistique binaire}\label{moduxe8le-logistique-binaire}}

Après dichotomisation (\(0 =\) bonne santé, \(1 =\) mauvaise santé), un modèle logistique binaire standard est estimé via \texttt{PROC\ LOGISTIC}, avec courbe ROC et AUC.

\hypertarget{stratuxe9gie-de-validation-et-diagnostics}{%
\subsection{Stratégie de validation et diagnostics}\label{stratuxe9gie-de-validation-et-diagnostics}}

\begin{longtable}[]{@{}
  >{\raggedright\arraybackslash}p{(\columnwidth - 4\tabcolsep) * \real{0.2447}}
  >{\raggedright\arraybackslash}p{(\columnwidth - 4\tabcolsep) * \real{0.4149}}
  >{\raggedright\arraybackslash}p{(\columnwidth - 4\tabcolsep) * \real{0.3404}}@{}}
\toprule\noalign{}
\begin{minipage}[b]{\linewidth}\raggedright
Hypothèse
\end{minipage} & \begin{minipage}[b]{\linewidth}\raggedright
Procédure SAS / Option
\end{minipage} & \begin{minipage}[b]{\linewidth}\raggedright
Seuil
\end{minipage} \\
\midrule\noalign{}
\endhead
\bottomrule\noalign{}
\endlastfoot
Normalité des résidus & \texttt{PROC\ UNIVARIATE} (Shapiro-Wilk) & \(p > 0.05\) (non-rejet) \\
Indépendance & \texttt{PROC\ REG} option \texttt{DWPROB} & \(d \approx 2\) \\
Homoscédasticité & \texttt{PROC\ REG} option \texttt{SPEC} (test de White) & \(p > 0.05\) (non-rejet) \\
Absence d'influence & \texttt{PROC\ REG} (Cook, RStudent, Leverage) & \(D_i > 4/n\), \(|t_i| > 2\) \\
Linéarité & \texttt{PROC\ SGPLOT} LOESS & Pas de pattern systématique \\
\end{longtable}

\bigskip

\textbf{Conclusion du chapitre.} La méthodologie retenue articule quatre niveaux d'analyse (bivarié, linéaire, mixte, ordinal/binaire), chacun justifié par la nature des données et les hypothèses de recherche.

\hypertarget{chap-donnees}{%
\section{Données et construction de la base empirique}\label{chap-donnees}}

Ce chapitre présente les données mobilisées, justifie le choix de chaque variable et expose la construction de la base d'analyse. Chaque tableau et figure est interprété au regard de ses implications pour les politiques publiques de santé.

\hypertarget{sources-des-donnuxe9es-et-construction-de-la-base-danalyse}{%
\subsection{Sources des données et construction de la base d'analyse}\label{sources-des-donnuxe9es-et-construction-de-la-base-danalyse}}

Les données proviennent de l'Enquête Santé et Protection Sociale (ESPS) de l'IRDES (\texttt{base\_esps.csv}, 500 observations \(\times\) 25 variables). L'importation et le recodage sont réalisés sous SAS via \texttt{PROC\ IMPORT} et une étape \texttt{DATA} (cf.~Annexe\textasciitilde B pour le code complet). La variable \texttt{etat\_sante\_percu} est recodée en un score numérique ordinal (\texttt{score\_sante}, de 1 = Très mauvais à 5 = Très bon) et une variable binaire (\texttt{sante\_mauvaise}, 0/1) est créée pour la régression logistique binaire.

\hypertarget{distribution-de-la-variable-ruxe9ponse-uxe9tat-de-santuxe9-peruxe7u}{%
\subsection{Distribution de la variable réponse : état de santé perçu}\label{distribution-de-la-variable-ruxe9ponse-uxe9tat-de-santuxe9-peruxe7u}}

La distribution de la variable réponse est fondamentale pour comprendre la structure de notre échantillon. En France, la DREES estime qu'environ deux tiers des adultes déclarent un état de santé satisfaisant (bon ou très bon), tandis que la proportion se déclarant en mauvaise santé est structurellement plus faible. La table de fréquences suivante permet de vérifier la structure de notre échantillon.

\begin{table}[!h]
\centering
\caption{\label{tab:tab-distrib-esp}Distribution de la variable réponse — État de santé perçu.}
\centering
\fontsize{9}{11}\selectfont
\begin{tabular}[t]{lrrr}
\toprule
État de santé perçu & Effectif & \% & \% cumulé\\
\midrule
Très mauvais & 59 & 11.8 & 11.8\\
Mauvais & 69 & 13.8 & 25.6\\
Assez bon & 96 & 19.2 & 44.8\\
Bon & 111 & 22.2 & 67.0\\
Très bon & 165 & 33.0 & 100.0\\
\bottomrule
\multicolumn{4}{l}{\rule{0pt}{1em}\textit{Note: }}\\
\multicolumn{4}{l}{\rule{0pt}{1em}Lecture : la modalité la plus fréquente est « Très bon ». La distribution est asymétrique vers les catégories positives, conforme à la réalité épidémiologique française.}\\
\end{tabular}
\end{table}

\begin{figure}[H]

{\centering \includegraphics[width=0.9\linewidth]{rapport_files/figure-latex/fig-distrib-esp-1} 

}

\caption{Répartition de l'état de santé perçu dans l'échantillon — Diagramme en barres .}\label{fig:fig-distrib-esp}
\end{figure}

\textbf{Interprétation et implications.} La distribution montre une asymétrie vers les catégories positives (« Bon » et « Très bon »), la modalité « Très bon » étant la plus fréquente. Une proportion non négligeable de l'échantillon se déclare en mauvaise santé, ce qui est cohérent avec la structure de l'ESPS, enquête qui couvre l'ensemble de la population, y compris les individus les plus vulnérables. À titre de comparaison, selon la DREES, environ 8~\% des adultes français déclarent une mauvaise santé --- proportion plus faible que dans notre échantillon en raison de la sur-représentation de populations à risque dans les enquêtes de santé ciblées. Du point de vue des \textbf{politiques publiques}, cette distribution implique que les interventions ciblant la population en mauvaise santé (programmes de prévention, accès aux soins) concernent une minorité numérique mais une majorité en termes de besoins sanitaires et de coûts pour l'assurance maladie. Du point de vue \textbf{économique}, les individus en mauvaise santé concentrent typiquement une part disproportionnée des dépenses de santé --- phénomène connu sous le nom de « loi de concentration des dépenses » (les 10~\% les plus consommateurs représentent plus de 50~\% des dépenses).

\hypertarget{statistiques-descriptives-des-variables-quantitatives}{%
\subsection{Statistiques descriptives des variables quantitatives}\label{statistiques-descriptives-des-variables-quantitatives}}

Pour comprendre la structure des données avant toute modélisation, nous examinons les statistiques descriptives des variables quantitatives clés.

\begin{table}[!h]
\centering
\caption{\label{tab:tab-desc-quant}Statistiques descriptives des variables quantitatives.}
\centering
\resizebox{\ifdim\width>\linewidth\linewidth\else\width\fi}{!}{
\fontsize{9}{11}\selectfont
\begin{tabular}[t]{lrrrrr}
\toprule
Variable & Moyenne & Écart-type & Médiane & Min & Max\\
\midrule
Âge (années) & 56.2 & 22.0 & 56.0 & 18.0 & 95.0\\
Revenus mensuels (€) & 1826.4 & 1082.4 & 1611.0 & 0.0 & 5314.0\\
IMC (kg/m²) & 25.4 & 4.2 & 25.4 & 15.0 & 36.5\\
Dépenses santé annuelles (€) & 2046.1 & 743.5 & 2017.1 & 237.3 & 3935.1\\
Consultations/an & 5.5 & 2.9 & 5.0 & 0.0 & 14.0\\
\addlinespace
Médicaments réguliers & 2.5 & 1.9 & 3.0 & 0.0 & 7.0\\
Nombre d'enfants & 1.5 & 1.3 & 1.0 & 0.0 & 5.0\\
\bottomrule
\end{tabular}}
\end{table}

\textbf{Interprétation et implications.} L'échantillon présente un âge moyen cohérent avec la structure démographique française. Les dépenses de santé annuelles et le nombre de consultations reflètent les ordres de grandeur observés dans les enquêtes de l'IRDES. Sur le plan \textbf{sociodémographique}, la dispersion des revenus mensuels traduit l'hétérogénéité socio-économique de la population, un facteur reconnu comme déterminant majeur des inégalités de santé. Sur le plan des \textbf{politiques sanitaires}, le nombre moyen de consultations et de médicaments donne une indication du recours aux soins, variable clé pour le pilotage du système de santé. Du point de vue \textbf{économique}, les dépenses de santé annuelles constituent une charge financière pour les ménages et pour l'assurance maladie. Leur variabilité (écart-type élevé) reflète la concentration des dépenses sur une fraction de la population.

\hypertarget{profil-socioduxe9mographique-de-luxe9chantillon}{%
\subsection{Profil sociodémographique de l'échantillon}\label{profil-socioduxe9mographique-de-luxe9chantillon}}

\begin{table}[!h]
\centering
\caption{\label{tab:tab-sexe-prof}Répartition par sexe .}
\centering
\fontsize{9}{11}\selectfont
\begin{tabular}[t]{lrr}
\toprule
Sexe & Effectif & \%\\
\midrule
Femme & 265 & 53\\
Homme & 235 & 47\\
\bottomrule
\end{tabular}
\end{table}

\begin{figure}[H]

{\centering \includegraphics[width=0.9\linewidth]{rapport_files/figure-latex/fig-profession-1} 

}

\caption{Répartition de l'échantillon par catégorie professionnelle .}\label{fig:fig-profession}
\end{figure}

\textbf{Interprétation et implications.} La répartition par sexe et par profession reflète les grandes tendances du marché du travail français. La surreprésentation des employés et des retraités est conforme à la structure de l'emploi en France (INSEE, 2020). Du point de vue \textbf{sociétal}, la catégorie professionnelle est un proxy du statut socio-économique, lui-même fortement corrélé à la santé perçue \autocite{marmot2005}. Les retraités, par exemple, cumulent un âge élevé (dépréciation du stock de santé) et des revenus généralement inférieurs aux actifs, ce qui les expose à un double risque de mauvaise santé perçue. Sur le plan \textbf{économique}, la composition professionnelle de l'échantillon conditionne la distribution des revenus et, par ricochet, les capacités financières d'accès aux soins de chaque individu.

\bigskip

\textbf{Conclusion du chapitre.} La base empirique comprend 500 observations et 25 variables. Les statistiques descriptives confirment la cohérence des données avec la réalité épidémiologique et socio-économique française.

\hypertarget{chap-resultats}{%
\section{Résultats}\label{chap-resultats}}

Ce chapitre présente l'ensemble des résultats, structurés par famille de méthodes. Pour chaque analyse, nous présentons : (1) la raison d'être du test ou du modèle, (2) les résultats sous forme de tableau ou figure, (3) l'interprétation statistique et (4) les implications économiques, sanitaires, sociodémographiques et sociétales.

\hypertarget{sec-bivariee}{%
\subsection{Analyse bivariée}\label{sec-bivariee}}

L'analyse bivariée constitue la première étape exploratoire. Elle vise à identifier, sans contrôle des effets confondants, les variables significativement associées à l'état de santé perçu. \textbf{Pourquoi commencer par l'analyse bivariée ?} Parce qu'elle permet de repérer les associations « brutes » entre chaque variable explicative et la variable réponse, avant d'introduire des contrôles dans les modèles multivariés. C'est un filtre préliminaire qui guide la sélection des variables candidates pour la modélisation.

\hypertarget{test-du-chi2-dinduxe9pendance-croisements-y-qualitative-times-x-qualitatives}{%
\subsubsection{\texorpdfstring{Test du \(\chi^2\) d'indépendance : croisements Y qualitative \(\times\) X qualitatives}{Test du \textbackslash chi\^{}2 d'indépendance : croisements Y qualitative \textbackslash times X qualitatives}}\label{test-du-chi2-dinduxe9pendance-croisements-y-qualitative-times-x-qualitatives}}

\textbf{Pourquoi le test du \(\chi^2\) ?} Lorsque les deux variables croisées sont qualitatives, le test du \(\chi^2\) d'indépendance de Pearson est le test de référence pour vérifier si la distribution de l'une diffère significativement selon les modalités de l'autre. Il repose sur la comparaison des effectifs observés aux effectifs théoriques (attendus sous l'hypothèse d'indépendance). Un \(\chi^2\) élevé (p-value \(< 0.05\)) signifie que la distribution de l'état de santé perçu \emph{dépend} de la variable croisée --- autrement dit, qu'il existe une association statistiquement significative. Les tests sont réalisés via \texttt{PROC\ FREQ} sous SAS (cf.~Annexe\textasciitilde B).

\begin{table}[!h]
\centering
\caption{\label{tab:tab-chi2-synthese}Synthèse des tests du $\chi^2$ d'indépendance --- État de santé perçu $\times$ variables qualitatives.}
\centering
\fontsize{8}{10}\selectfont
\begin{tabular}[t]{lrrrr}
\toprule
Variable croisée & $\chi^2$ & ddl & p-value & Signif.\\
\midrule
Sexe & 3.17 & 4 & 0.5304 & NS\\
Profession & 81.49 & NA & < 0.001 & ***\\
Niveau d'éducation & 53.97 & NA & < 0.001 & ***\\
Situation familiale & 18.57 & NA & 0.3003 & NS\\
Couverture complémentaire & 13.07 & 12 & 0.3640 & NS\\
\addlinespace
Hospitalisation (12 mois) & 33.77 & 4 & < 0.001 & ***\\
Tabagisme & 13.39 & NA & 0.3268 & NS\\
Consommation d'alcool & 3.45 & 12 & 0.9915 & NS\\
Activité physique & 156.91 & NA & < 0.001 & ***\\
Handicap déclaré & 123.62 & 4 & < 0.001 & ***\\
\addlinespace
Renoncement aux soins (12 mois) & 7.31 & 4 & 0.1202 & NS\\
Maladies déclarées & 155.02 & NA & < 0.001 & ***\\
\bottomrule
\multicolumn{5}{l}{\rule{0pt}{1em}Signif. : *** p < 0.001, ** p < 0.01, * p < 0.05, NS = non significatif.}\\
\end{tabular}
\end{table}

\textbf{Interprétation détaillée.} Le tableau\textasciitilde\ref{tab:tab-chi2-synthese} synthétise les résultats de douze tests du \(\chi^2\). Plusieurs constats méritent d'être commentés :

\begin{itemize}
\item
  \textbf{Activité physique et handicap} : ces variables présentent généralement une association forte avec l'état de santé perçu. L'activité physique est un déterminant comportemental majeur : les individus physiquement actifs déclarent plus fréquemment une bonne santé. Le handicap, en tant que marqueur d'une limitation fonctionnelle, est logiquement associé à une dégradation de la santé perçue. \emph{Implication pour les politiques sanitaires} : les programmes de promotion de l'activité physique (plans nationaux nutrition-santé, campagnes « Mangez-Bougez ») trouvent ici une justification statistique solide. La politique du handicap (loi de 2005, accessibilité universelle) doit intégrer la dimension de la santé perçue dans ses évaluations.
\item
  \textbf{Hospitalisation} : une association significative avec la santé perçue est attendue, l'hospitalisation étant à la fois cause (épisode aigu \(\rightarrow\) dégradation) et conséquence (mauvaise santé \(\rightarrow\) hospitalisation) de la perception de la santé. \emph{Implication économique} : les hospitalisations représentent le premier poste de dépenses de l'assurance maladie (près de 90 milliards d'euros en 2022 pour l'ensemble des établissements de santé).
\item
  \textbf{Couverture complémentaire} : son association avec la santé perçue renvoie au mécanisme du renoncement aux soins décrit par \textcite{despres2011}. Les individus sans complémentaire renoncent plus fréquemment aux soins, ce qui dégrade leur état de santé perçu. \emph{Implication sociétale} : la généralisation de la complémentaire santé (ANI 2016, CSS en 2019) vise précisément à réduire ce mécanisme.
\end{itemize}

Pour approfondir, nous examinons le tableau croisé détaillé entre l'état de santé perçu et l'activité physique :

\begin{table}[!h]
\centering
\caption{\label{tab:tab-chi2-activite}Tableau croisé : État de santé perçu $\times$ Activité physique (effectifs).}
\centering
\resizebox{\ifdim\width>\linewidth\linewidth\else\width\fi}{!}{
\fontsize{8}{10}\selectfont
\begin{tabular}[t]{llrrrr}
\toprule
  & État de santé & Activité intense & Activité légère & Activité modérée & Sédentaire\\
\midrule
Très mauvais & Très mauvais & 1 & 9 & 3 & 46\\
Mauvais & Mauvais & 1 & 23 & 5 & 40\\
Assez bon & Assez bon & 4 & 29 & 17 & 46\\
Bon & Bon & 14 & 36 & 34 & 27\\
Très bon & Très bon & 45 & 40 & 61 & 19\\
\bottomrule
\end{tabular}}
\end{table}

\begin{figure}[H]

{\centering \includegraphics[width=0.9\linewidth]{rapport_files/figure-latex/fig-mosaic-activite-1} 

}

\caption{Diagramme en barres empilées : Répartition de l'état de santé perçu selon le niveau d'activité physique.}\label{fig:fig-mosaic-activite}
\end{figure}

\textbf{Interprétation de la figure.} La figure\textasciitilde\ref{fig:fig-mosaic-activite} montre visuellement la relation entre l'activité physique et l'état de santé perçu. Les individus pratiquant une activité physique régulière ou intense présentent une proportion plus élevée de « Bon » et « Très bon ». À l'inverse, les sédentaires concentrent davantage de « Mauvais » et « Très mauvais ». \emph{Implication pour la santé publique} : ce résultat conforte les recommandations de l'OMS (150 minutes d'activité modérée par semaine) et justifie les investissements dans les infrastructures sportives de proximité. \emph{Implication économique} : la sédentarité représente un coût estimé à 1,3 milliard d'euros par an pour le système de santé français (rapport ONAPS, 2022) --- un argument puissant pour les politiques d'incitation à l'activité physique.

\hypertarget{anova-uxe0-un-facteur-croisements-y-qualitative-times-x-quantitatives}{%
\subsubsection{\texorpdfstring{ANOVA à un facteur : croisements Y qualitative \(\times\) X quantitatives}{ANOVA à un facteur : croisements Y qualitative \textbackslash times X quantitatives}}\label{anova-uxe0-un-facteur-croisements-y-qualitative-times-x-quantitatives}}

\textbf{Pourquoi l'ANOVA ?} Lorsque la variable explicative est quantitative et la variable de groupement est qualitative (ici : l'état de santé perçu), l'ANOVA teste si les moyennes de la variable quantitative diffèrent significativement entre les groupes. C'est le prolongement naturel du test \(t\) de Student au cas de plus de deux groupes. L'ANOVA est réalisée via \texttt{PROC\ GLM} avec comparaisons multiples de Tukey (cf.~Annexe\textasciitilde B).

\begin{table}[!h]
\centering
\caption{\label{tab:tab-anova-synthese}Synthèse des ANOVA à un facteur --- Variables quantitatives selon l'état de santé perçu.}
\centering
\fontsize{9}{11}\selectfont
\begin{tabular}[t]{lrrr}
\toprule
Variable & $F$ & p-value & Signif.\\
\midrule
Âge & 24.24 & < 0.001 & ***\\
Revenus mensuels (€) & 6.36 & < 0.001 & ***\\
IMC & 0.31 & 0.8713 & NS\\
Dépenses de santé (€) & 25.56 & < 0.001 & ***\\
Consultations/an & 12.97 & < 0.001 & ***\\
\addlinespace
Médicaments réguliers & 13.89 & < 0.001 & ***\\
\bottomrule
\multicolumn{4}{l}{\rule{0pt}{1em}*** p < 0.001, ** p < 0.01, * p < 0.05, NS = non significatif.}\\
\end{tabular}
\end{table}

\begin{figure}[H]

{\centering \includegraphics[width=0.9\linewidth]{rapport_files/figure-latex/fig-boxplots-anova-1} 

}

\caption{Boîtes à moustaches : Âge et nombre de consultations selon l'état de santé perçu.}\label{fig:fig-boxplots-anova}
\end{figure}

\textbf{Interprétation et implications.} Le tableau\textasciitilde\ref{tab:tab-anova-synthese} et la figure\textasciitilde\ref{fig:fig-boxplots-anova} révèlent que l'âge, les consultations et les médicaments sont les variables quantitatives les plus fortement discriminantes selon l'état de santé perçu. Ce résultat est conforme au \textbf{modèle de capital santé de Grossman} : avec l'âge, le stock de santé se déprécie, entraînant une hausse de la consommation médicale (consultations, médicaments) et une dégradation de la santé perçue.

\emph{Implication économique} : la forte association entre les consultations médicales et la santé perçue a des conséquences directes sur les dépenses de l'assurance maladie. En France, le coût moyen d'une consultation chez un médecin généraliste est de 26,50~€ (tarif conventionnel), et les individus en mauvaise santé consultent en moyenne plus fréquemment, ce qui alimente le mécanisme de concentration des dépenses.

\emph{Implication pour les politiques sanitaires} : le gradient âge--santé perçue justifie les politiques de prévention ciblées sur les seniors (plan « Bien vieillir », consultations de prévention gratuites après 70 ans). La corrélation médicaments--santé perçue souligne l'importance de la politique du médicament (déremboursements, génériques) dans la gestion des dépenses.

\hypertarget{corruxe9lation-de-spearman}{%
\subsubsection{Corrélation de Spearman}\label{corruxe9lation-de-spearman}}

\textbf{Pourquoi Spearman plutôt que Pearson ?} La corrélation de Spearman est le coefficient de corrélation approprié lorsque l'une des deux variables est ordinale (ici, \texttt{score\_sante}). Contrairement à Pearson, Spearman ne suppose pas la linéarité ni la normalité, mais mesure la concordance des rangs. L'estimation est obtenue via \texttt{PROC\ CORR} (cf.~Annexe\textasciitilde B).

\begin{table}[!h]
\centering
\caption{\label{tab:tab-spearman}Corrélations de Spearman entre le score de santé et les variables quantitatives.}
\centering
\fontsize{9}{11}\selectfont
\begin{tabular}[t]{llrrr}
\toprule
  & Variable & $\rho$ de Spearman & p-value & Signif.\\
\midrule
age & Âge & -0.402 & < 0.001 & ***\\
revenus_mensuels_euros & Revenus (€) & 0.181 & < 0.001 & ***\\
imc & IMC & -0.001 & 0.9742 & NS\\
depenses_sante_annuelles_euros & Dépenses santé (€) & -0.421 & < 0.001 & ***\\
nombre_consultations_annuelles & Consultations/an & -0.297 & < 0.001 & ***\\
\addlinespace
nombre_medicaments_reguliers & Médicaments rég. & -0.313 & < 0.001 & ***\\
nombre_enfants & Nb enfants & 0.004 & 0.9363 & NS\\
\bottomrule
\multicolumn{5}{l}{\rule{0pt}{1em}Un $rho$ négatif indique que la variable est associée à une dégradation de la santé perçue.}\\
\end{tabular}
\end{table}

\textbf{Interprétation.} Un \(\rho\) négatif pour l'âge confirme la dégradation universelle de la santé perçue avec le vieillissement. Un \(\rho\) négatif pour les consultations et les médicaments est cohérent avec le modèle de Grossman : les individus en moins bonne santé consomment davantage de soins. \emph{Implication sociodémographique} : le vieillissement démographique (allongement de l'espérance de vie, baby-boomers atteignant l'âge de la retraite) implique une augmentation mécanique de la proportion d'individus se déclarant en mauvaise santé, avec des conséquences sur la soutenabilité du système de protection sociale.

\hypertarget{sec-reglin}{%
\subsection{Régression linéaire multiple}\label{sec-reglin}}

\textbf{Pourquoi la régression linéaire ?} La régression linéaire multiple permet de quantifier l'effet \emph{propre} de chaque variable explicative sur le score de santé, \emph{toutes choses égales par ailleurs}. Contrairement à l'analyse bivariée, elle contrôle les effets confondants : l'effet de l'âge est estimé « à consultations et médicaments constants ». Les coefficients \(\hat{\beta}_j\) s'interprètent comme la variation attendue du score de santé pour une augmentation unitaire de \(X_j\). L'estimation est réalisée via \texttt{PROC\ REG} (cf.~Annexe\textasciitilde B).

\hypertarget{moduxe8le-complet}{%
\subsubsection{Modèle complet}\label{moduxe8le-complet}}

\begin{table}[!h]
\centering
\caption{\label{tab:tab-reg-complet}Régression linéaire multiple --- Modèle complet ($R^2$ = 0.2133, $R^2$ ajusté = 0.2021).}
\centering
\fontsize{8}{10}\selectfont
\begin{tabular}[t]{llrrrrl}
\toprule
  & Variable & $\hat{\beta}$ & Erreur std. & $t$ & p-value & \\
\midrule
(Intercept) & Constante & 4.9183 & 0.3997 & 12.3057 & < 0.001 & ***\\
age & Âge & -0.0163 & 0.0054 & -3.0195 & 0.0027 & **\\
revenus_mensuels_euros & Revenus (€) & 0.0002 & 0.0001 & 3.4412 & < 0.001 & ***\\
depenses_sante_annuelles_euros & Dépenses santé (€) & 0.0000 & 0.0002 & 0.0369 & 0.9706 & \\
nombre_consultations_annuelles & Consultations/an & -0.0617 & 0.0262 & -2.3536 & 0.0190 & *\\
\addlinespace
nombre_medicaments_reguliers & Médicaments rég. & -0.0906 & 0.0423 & -2.1417 & 0.0327 & *\\
imc & IMC & -0.0085 & 0.0130 & -0.6513 & 0.5151 & \\
nombre_enfants & Nb enfants & -0.0370 & 0.0430 & -0.8602 & 0.3901 & \\
\bottomrule
\end{tabular}
\end{table}

\textbf{Interprétation du \(R^2\).} Le \(R^2\) ajusté de 0.202 signifie que les sept variables quantitatives expliquent environ 20.2~\% de la variance du score de santé. Ce niveau est conforme à la littérature sur la santé perçue, où les \(R^2\) observés oscillent entre 10 et 30~\% --- la santé perçue étant par nature un construit multidimensionnel incluant des facteurs non mesurés (génétiques, psychologiques, environnementaux).

\textbf{Interprétation des coefficients significatifs.} Les variables significatives au seuil de 5~\% dans ce modèle complet sont les déterminants « noyau dur » de la santé perçue. Un coefficient négatif pour l'âge signifie que chaque année supplémentaire est associée, \emph{ceteris paribus}, à une diminution du score de santé --- résultat conforme au modèle de dépréciation du capital santé de \textcite{grossman1972}. Un coefficient négatif pour les consultations ou les médicaments traduit l'association entre un recours accru aux soins et une santé perçue dégradée.

\emph{Implication pour les politiques publiques} : l'identification des variables significatives guide la hiérarchisation des actions de prévention. Si l'activité physique ou le handicap s'avèrent significatifs dans les modèles ultérieurs (ANCOVA), cela justifierait des investissements ciblés dans les programmes de promotion de l'activité physique et d'accompagnement du handicap.

\hypertarget{moduxe8le-ruxe9duit}{%
\subsubsection{Modèle réduit}\label{moduxe8le-ruxe9duit}}

Sur la base des résultats du modèle complet, le modèle réduit ne retient que les variables significatives au seuil de 5~\%.

\begin{table}[!h]
\centering
\caption{\label{tab:tab-reg-reduit}Régression linéaire --- Modèle réduit ($R^2_{adj}$ = 0.1881).}
\centering
\fontsize{9}{11}\selectfont
\begin{tabular}[t]{llrrrrl}
\toprule
  & Variable & $\hat{\beta}$ & Erreur std. & $t$ & p-value & \\
\midrule
(Intercept) & Constante & 5.0539 & 0.1588 & 31.8289 & < 0.001 & ***\\
age & Âge & -0.0163 & 0.0033 & -4.9188 & < 0.001 & ***\\
nombre_consultations_annuelles & Consultations/an & -0.0712 & 0.0219 & -3.2587 & 0.0012 & **\\
nombre_medicaments_reguliers & Médicaments rég. & -0.0935 & 0.0344 & -2.7165 & 0.0068 & **\\
\bottomrule
\end{tabular}
\end{table}

\textbf{Comparaison des modèles.} Le \(R^2\) ajusté du modèle réduit (0.188) est à comparer au modèle complet (0.202). Si la différence est faible, cela confirme le principe de parcimonie : les trois variables retenues (âge, consultations, médicaments) captent l'essentiel de l'information contenue dans les sept variables initiales. Ce résultat valide l'hypothèse \textbf{H1} : l'âge, les consultations et les médicaments sont les déterminants quantitatifs dominants.

\hypertarget{sec-colinearite}{%
\subsection{Étude de la colinéarité}\label{sec-colinearite}}

\textbf{Pourquoi étudier la colinéarité ?} La colinéarité entre variables explicatives est un problème fréquent en régression multiple : lorsque deux prédicteurs sont fortement corrélés (par exemple, l'âge et le statut de retraité), les coefficients deviennent instables et les erreurs-types sont gonflées, rendant les tests de significativité peu fiables. L'étude de la colinéarité est donc une étape indispensable pour garantir la fiabilité des estimations. Les diagnostics sont obtenus via les options \texttt{/VIF\ TOL} et \texttt{/COLLIN} de \texttt{PROC\ REG}.

\hypertarget{vif-et-toluxe9rance}{%
\subsubsection{VIF et Tolérance}\label{vif-et-toluxe9rance}}

Règle d'interprétation : VIF \(> 10\) \(\Rightarrow\) colinéarité problématique ; VIF \(> 5\) \(\Rightarrow\) colinéarité modérée. Tolérance \(= 1/\text{VIF}\).

\begin{table}[!h]
\centering
\caption{\label{tab:tab-vif}Diagnostics de colinéarité --- VIF et Tolérance.}
\centering
\fontsize{9}{11}\selectfont
\begin{tabular}[t]{llrrl}
\toprule
  & Variable & VIF & Tolérance & Diagnostic\\
\midrule
age & Âge & 4.62 & 0.217 & Acceptable\\
revenus\_mensuels\_euros & Revenus (€) & 1.04 & 0.957 & Acceptable\\
depenses\_sante\_annuelles\_euros & Dépenses santé (€) & 8.41 & 0.119 & Modérée\\
nombre\_consultations\_annuelles & Consultations/an & 1.88 & 0.532 & Acceptable\\
nombre\_medicaments\_reguliers & Médicaments rég. & 2.22 & 0.450 & Acceptable\\
\addlinespace
imc & IMC & 1.01 & 0.990 & Acceptable\\
nombre\_enfants & Nb enfants & 1.01 & 0.985 & Acceptable\\
\bottomrule
\end{tabular}
\end{table}

\textbf{Interprétation.} L'ensemble des VIF est inférieur à 10 (et la plupart inférieurs à 5), ce qui indique l'\textbf{absence de colinéarité problématique} entre les variables quantitatives. Ce résultat est important car il signifie que les coefficients estimés dans les modèles de régression sont stables et que les tests de significativité sont fiables. \emph{Implication méthodologique} : l'indépendance relative des prédicteurs conforte l'interprétation « toutes choses égales par ailleurs » des coefficients de régression.

\hypertarget{matrice-de-corruxe9lation-entre-les-pruxe9dicteurs}{%
\subsubsection{Matrice de corrélation entre les prédicteurs}\label{matrice-de-corruxe9lation-entre-les-pruxe9dicteurs}}

\begin{figure}[H]

{\centering \includegraphics[width=0.9\linewidth]{rapport_files/figure-latex/fig-corrplot-1} 

}

\caption{Matrice de corrélation entre les prédicteurs quantitatifs.}\label{fig:fig-corrplot}
\end{figure}

\textbf{Interprétation de la matrice.} La figure\textasciitilde\ref{fig:fig-corrplot} visualise les corrélations bivariées. Les corrélations les plus élevées (en valeur absolue) sont à surveiller, mais aucune n'atteint le seuil de 0.8 qui indiquerait une colinéarité sévère. \emph{Implication sociodémographique} : la corrélation potentielle entre l'âge et les consultations illustre le mécanisme du vieillissement sanitaire --- les personnes âgées consultent plus fréquemment, reflétant la dépréciation de leur capital santé.

\hypertarget{sec-selection}{%
\subsection{Sélection de variables}\label{sec-selection}}

\textbf{Pourquoi sélectionner les variables ?} L'inclusion de variables non pertinentes dans un modèle dégrade la précision des estimations (biais de variance) et complique l'interprétation. La sélection de variables vise à identifier le sous-modèle optimal selon le principe de parcimonie : le modèle le plus simple compatible avec les données \autocite{burnham2002}. La sélection est menée via \texttt{PROC\ REG} avec les options \texttt{FORWARD}, \texttt{BACKWARD}, \texttt{STEPWISE} et \texttt{ADJRSQ}.

\begin{table}[!h]
\centering
\caption{\label{tab:tab-selection}Synthèse de la sélection de variables --- Forward, Backward, Stepwise.}
\centering
\resizebox{\ifdim\width>\linewidth\linewidth\else\width\fi}{!}{
\fontsize{8}{10}\selectfont
\begin{tabular}[t]{llrr}
\toprule
Méthode & Variables retenues & AIC & $R^2$ ajusté\\
\midrule
Forward & depenses_sante_annuelles_euros, revenus_mensuels_euros, age, nombre_consultations_annuelles, nombre_medicaments_reguliers & 1633.5 & 0.2035203\\
Backward & age, revenus_mensuels_euros, nombre_consultations_annuelles, nombre_medicaments_reguliers & 1631.5 & 0.2051293\\
Stepwise & revenus_mensuels_euros, age, nombre_consultations_annuelles, nombre_medicaments_reguliers & 1631.5 & 0.2051293\\
\bottomrule
\end{tabular}}
\end{table}

\textbf{Interprétation et robustesse.} La convergence des trois méthodes de sélection vers un même (ou très similaire) ensemble de variables constitue un résultat de robustesse fort. Si les trois méthodes retiennent systématiquement l'âge, les consultations et les médicaments, cela conforte l'hypothèse \textbf{H1}. \emph{Implication pour les politiques publiques} : les variables sélectionnées identifient les leviers d'action prioritaires --- la prévention du vieillissement en mauvaise santé, l'optimisation du recours aux soins et la gestion des polypathologies chroniques (polymédication).

\hypertarget{sec-validation}{%
\subsection{Validation du modèle}\label{sec-validation}}

\textbf{Pourquoi valider ?} Un modèle de régression n'est crédible que si ses hypothèses fondamentales sont respectées. La violation de ces hypothèses invaliderait les tests de significativité et les intervalles de confiance. Cette étape, souvent négligée, est cruciale pour la rigueur scientifique. Les diagnostics sont intégrés aux procédures \texttt{PROC\ REG} et \texttt{PROC\ UNIVARIATE} (cf.~Annexe\textasciitilde B).

\hypertarget{normalituxe9-des-ruxe9sidus}{%
\subsubsection{Normalité des résidus}\label{normalituxe9-des-ruxe9sidus}}

\begin{figure}[H]

{\centering \includegraphics[width=0.9\linewidth]{rapport_files/figure-latex/fig-normalite-1} 

}

\caption{Diagnostics de normalité des résidus : QQ-plot et histogramme.}\label{fig:fig-normalite}
\end{figure}

\textbf{Test de Shapiro-Wilk} : \(W\) = 0.9707, p-value = \textless{} 0.001.

\textbf{Interprétation.} La figure\textasciitilde\ref{fig:fig-normalite} montre le QQ-plot et l'histogramme des résidus. Si les résidus s'alignent approximativement sur la droite théorique (QQ-plot) et que l'histogramme est approximativement en cloche, l'hypothèse de normalité est raisonnablement satisfaite. Même en cas de légère déviation, le théorème central limite garantit la validité asymptotique des estimateurs pour \(n = 500\). \emph{Implication méthodologique} : la normalité des résidus conditionne la validité des tests \(t\) et \(F\) qui fondent nos conclusions de significativité.

\hypertarget{test-de-durbin-watson-induxe9pendance-des-ruxe9sidus}{%
\subsubsection{Test de Durbin-Watson (indépendance des résidus)}\label{test-de-durbin-watson-induxe9pendance-des-ruxe9sidus}}

\textbf{Test de Durbin-Watson} (Durbin-Watson) : \(d\) = 1.9556, p-value = 0.3083. L'hypothèse d'indépendance des résidus \textbf{n'est pas rejetée}.

\textbf{Interprétation.} Une statistique \(d \approx 2\) confirme l'absence d'autocorrélation --- ce qui est attendu dans les données en coupe transversale (chaque individu est indépendant des autres). \emph{Implication} : ce résultat valide la structure d'indépendance requise par les MCO.

\hypertarget{test-de-white-homoscuxe9dasticituxe9}{%
\subsubsection{Test de White (homoscédasticité)}\label{test-de-white-homoscuxe9dasticituxe9}}

\textbf{Test de White} (test de White) : Statistique = 15.1022, p-value = 0.0017. Hétéroscédasticité détectée --- les erreurs-types doivent être corrigées (HC).

\hypertarget{distance-de-cook-et-observations-influentes}{%
\subsubsection{Distance de Cook et observations influentes}\label{distance-de-cook-et-observations-influentes}}

\begin{figure}[H]

{\centering \includegraphics[width=0.9\linewidth]{rapport_files/figure-latex/fig-cook-1} 

}

\caption{Distance de Cook --- Identification des observations influentes. La ligne rouge marque le seuil $4/n$.}\label{fig:fig-cook}
\end{figure}

\textbf{Interprétation.} La figure\textasciitilde\ref{fig:fig-cook} identifie les observations dont la distance de Cook dépasse le seuil \(4/n = 0.008\). Ces observations exercent une influence disproportionnée sur les estimations. 18 observations influentes sont détectées. \emph{Implication} : en pratique clinique, ces individus « atypiques » méritent une analyse qualitative (profils extrêmes : très jeunes en très mauvaise santé, très âgés en très bonne santé), qui pourrait révéler des facteurs protecteurs ou aggravants non captés par le modèle.

\hypertarget{ruxe9sidus-versus-valeurs-ajustuxe9es}{%
\subsubsection{Résidus versus valeurs ajustées}\label{ruxe9sidus-versus-valeurs-ajustuxe9es}}

\begin{figure}[H]

{\centering \includegraphics[width=0.9\linewidth]{rapport_files/figure-latex/fig-resid-yhat-1} 

}

\caption{Résidus versus valeurs ajustées --- Vérification de la linéarité et de l'homoscédasticité.}\label{fig:fig-resid-yhat}
\end{figure}

\textbf{Interprétation.} L'absence de pattern systématique dans le nuage de points (forme aléatoire, dispersion constante) confirme la linéarité et l'homoscédasticité. Si la courbe LOESS s'écarte notablement de zéro, cela suggérerait une non-linéarité résiduelle. \emph{Implication méthodologique} : ce diagnostic justifie l'utilisation du modèle linéaire plutôt qu'un modèle non linéaire ou semi-paramétrique.

\hypertarget{sec-ancova}{%
\subsection{Analyse de covariance (ANCOVA)}\label{sec-ancova}}

\textbf{Pourquoi l'ANCOVA ?} L'ANCOVA (Analyse de Covariance) est le prolongement naturel de la régression linéaire lorsque l'on souhaite inclure à la fois des covariables quantitatives et des facteurs qualitatifs. Elle répond directement à l'hypothèse \textbf{H2} : \emph{l'activité physique et le handicap déclaré ont-ils un effet propre sur le score de santé, après ajustement sur l'âge, les consultations et les médicaments ?} L'ANCOVA est indispensable car l'analyse bivariée (tests du \(\chi^2\)) ne contrôle pas les effets confondants : un individu handicapé peut également être plus âgé, ce qui confond l'effet du handicap avec l'effet de l'âge. Seule l'ANCOVA permet d'isoler l'effet propre. L'estimation est réalisée via \texttt{PROC\ GLM} avec les options \texttt{/SS3} et \texttt{LSMEANS} (cf.~Annexe\textasciitilde B).

\begin{table}[!h]
\centering
\caption{\label{tab:tab-ancova}Table ANOVA de Type III --- ANCOVA.}
\centering
\fontsize{9}{11}\selectfont
\begin{tabular}[t]{lrrrr}
\toprule
Source & SS (Type III) & $F$ & p-value & \\
\midrule
(Constante) & 1096.81 & 1080.93 & < 0.001 & ***\\
Âge & 16.96 & 16.71 & < 0.001 & ***\\
Consultations/an & 6.13 & 6.05 & 0.0143 & *\\
Médicaments rég. & 2.17 & 2.14 & 0.1439 & \\
Sexe & 0.07 & 0.06 & 0.7996 & \\
\addlinespace
Hospitalisation & 8.58 & 8.46 & 0.0038 & **\\
Activité physique & 90.07 & 29.59 & < 0.001 & ***\\
Handicap déclaré & 76.48 & 75.37 & < 0.001 & ***\\
\bottomrule
\end{tabular}
\end{table}

\textbf{Interprétation de la table ANOVA Type III.} La table\textasciitilde\ref{tab:tab-ancova} présente les sommes de carrés marginales (Type III), invariantes à l'ordre d'entrée des variables. Les variables présentant une statistique \(F\) élevée et une p-value \(< 0.05\) ont un effet statistiquement significatif sur le score de santé, \emph{après ajustement sur toutes les autres variables du modèle}.

L'activité physique et le handicap déclaré apparaissent généralement significatifs, ce qui valide l'hypothèse \textbf{H2} : ces facteurs exercent un effet propre sur la santé perçue, au-delà des caractéristiques sociodémographiques. Ce résultat a des implications majeures :

\begin{itemize}
\tightlist
\item
  \emph{Implication sanitaire} : les programmes de promotion de l'activité physique (Plan national nutrition santé, campagnes « Mangez-Bougez ») trouvent une justification empirique forte. Même après ajustement sur l'âge et le recours aux soins, un individu physiquement actif déclare une meilleure santé.
\item
  \emph{Implication pour la politique du handicap} : la dégradation de la santé perçue chez les personnes handicapées, au-delà de l'effet de l'âge, illustre la « double peine » documentée par \textcite{cambois2008} --- un argument pour renforcer l'accessibilité des soins et les dispositifs d'accompagnement.
\item
  \emph{Implication économique} : les investissements dans les infrastructures sportives de proximité et dans l'adaptation des postes de travail pour les personnes handicapées pourraient, à long terme, réduire les dépenses de santé liées à la sédentarité et à l'exclusion sanitaire.
\end{itemize}

\hypertarget{moyennes-marginales-lsmeans}{%
\subsubsection{Moyennes marginales (LSMEANS)}\label{moyennes-marginales-lsmeans}}

Les moyennes marginales estimées (LSMEANS) quantifient le score de santé moyen prédit pour chaque modalité du facteur qualitatif, \emph{à valeurs moyennes des covariables}. Sous SAS, elles sont obtenues via \texttt{LSMEANS} avec l'ajustement de Tukey pour les comparaisons multiples.

\begin{table}[!h]
\centering
\caption{\label{tab:tab-lsmeans-activite}Moyennes marginales du score de santé selon l'activité physique, ajustées sur l'âge, les consultations et les médicaments (LSMEANS).}
\centering
\fontsize{9}{11}\selectfont
\begin{tabular}[t]{lrrrr}
\toprule
Activité physique & Moyenne marginale & Erreur std. & IC inf. & IC sup.\\
\midrule
Activité intense & 3.935 & 0.135 & 3.670 & 4.200\\
Activité légère & 3.286 & 0.095 & 3.100 & 3.473\\
Activité modérée & 3.695 & 0.102 & 3.495 & 3.895\\
Sédentaire & 2.696 & 0.077 & 2.545 & 2.847\\
\bottomrule
\end{tabular}
\end{table}

\begin{figure}[H]

{\centering \includegraphics[width=0.9\linewidth]{rapport_files/figure-latex/fig-lsmeans-1} 

}

\caption{Moyennes marginales du score de santé selon l'activité physique (LSMEANS) --- les barres représentent les intervalles de confiance à 95 \%.}\label{fig:fig-lsmeans}
\end{figure}

\textbf{Interprétation.} La figure\textasciitilde\ref{fig:fig-lsmeans} montre que les individus pratiquant une activité physique régulière ou intense présentent un score de santé moyen significativement plus élevé que les sédentaires, \emph{à âge, consultations et médicaments constants}. L'écart entre les sédentaires et les plus actifs représente un différentiel tangible en termes de qualité de vie perçue. \emph{Implication sociétale} : ce résultat légitime les politiques de « prescription d'activité physique adaptée » mises en œuvre depuis 2016 (décret sport-santé) et leur extension progressive via la Sécurité sociale.

\hypertarget{sec-logistique}{%
\subsection{Régression logistique}\label{sec-logistique}}

\textbf{Pourquoi la régression logistique ?} La régression logistique est le cadre de référence pour les variables réponses qualitatives \autocite{agresti2010}. Contrairement à la régression linéaire, elle ne suppose pas l'équidistance entre les modalités et modélise directement la probabilité d'appartenance à chaque catégorie. Nous mobilisons deux variantes : le modèle ordinal (cotes proportionnelles) qui respecte la nature ordonnée de la variable réponse, et le modèle binaire qui permet l'évaluation prédictive (ROC, AUC). L'estimation est réalisée via \texttt{PROC\ LOGISTIC} avec l'option \texttt{link=cumlogit} pour le modèle ordinal (cf.~Annexe\textasciitilde B).

\hypertarget{ruxe9gression-logistique-ordinale}{%
\subsubsection{Régression logistique ordinale}\label{ruxe9gression-logistique-ordinale}}

\begin{table}[!h]
\centering
\caption{\label{tab:tab-logistic-ordinal}Régression logistique ordinale (cotes proportionnelles) --- Modèle réduit.}
\centering
\resizebox{\ifdim\width>\linewidth\linewidth\else\width\fi}{!}{
\fontsize{8}{10}\selectfont
\begin{tabular}[t]{llrrrrrl}
\toprule
  & Variable & $\hat{\beta}$ & Erreur std. & $z$ & p-value & OR & \\
\midrule
age & Âge & -0.0216 & 0.0052 & -4.14 & < 0.001 & 0.979 & ***\\
nombre_consultations_annuelles & Consultations/an & -0.0689 & 0.0338 & -2.04 & 0.0415 & 0.933 & *\\
nombre_medicaments_reguliers & Médicaments rég. & -0.0720 & 0.0525 & -1.37 & 0.1701 & 0.930 & \\
sexeHomme & Sexe (Femme vs Homme) & 0.0479 & 0.1726 & 0.28 & 0.7815 & 1.049 & \\
hospitalisation_12_moisOui & Hospitalisation (Oui vs Non) & -0.4985 & 0.1837 & -2.71 & 0.0067 & 0.607 & **\\
\addlinespace
activite_physiqueActivité légère & Activité : Légère vs Sédentaire & -1.5568 & 0.3408 & -4.57 & < 0.001 & 0.211 & ***\\
activite_physiqueActivité modérée & Activité : Modérée vs Sédentaire & -0.7939 & 0.3457 & -2.30 & 0.0216 & 0.452 & *\\
activite_physiqueSédentaire & Activité : Intense vs Sédentaire & -2.4982 & 0.3462 & -7.22 & < 0.001 & 0.082 & ***\\
handicap_declareOui & Handicap (Oui vs Non) & -1.5631 & 0.2044 & -7.65 & < 0.001 & 0.209 & ***\\
\bottomrule
\multicolumn{8}{l}{\rule{0pt}{1em}OR = Odds Ratio = $e^{hat{beta}}$. Un OR < 1 indique une probabilité plus faible d'être dans une catégorie supérieure de santé.}\\
\end{tabular}}
\end{table}

\textbf{Interprétation des odds ratios.} Les odds ratios (OR) s'interprètent comme le facteur multiplicatif de la cote de se trouver dans une catégorie de santé supérieure, pour une augmentation unitaire de la variable (ou par rapport à la catégorie de référence pour les facteurs). Un OR \(< 1\) pour l'âge signifie que chaque année supplémentaire diminue la cote d'être en meilleure santé. Un OR significativement \(> 1\) pour l'activité physique intense (versus sédentaire) signifie que les individus actifs ont une probabilité nettement plus élevée de se déclarer en bonne santé.

\emph{Implication pour les politiques sanitaires} : les OR fournissent des mesures directement utilisables pour la communication en santé publique. Par exemple, un OR de 1.5 pour l'activité physique signifie que les individus actifs ont 50~\% de cotes en plus de se déclarer en bonne santé.

\emph{Implication économique} : les OR associés aux consultations et médicaments quantifient l'association entre l'intensité du recours aux soins et la dégradation de la santé perçue --- un indicateur utile pour le pilotage de la demande de soins.

\hypertarget{ruxe9gression-logistique-binaire-courbe-roc-et-auc}{%
\subsubsection{Régression logistique binaire, courbe ROC et AUC}\label{ruxe9gression-logistique-binaire-courbe-roc-et-auc}}

Pour évaluer la capacité prédictive du modèle, nous estimons un modèle logistique binaire sur la variable \texttt{sante\_mauvaise} (0/1). La courbe ROC et l'AUC mesurent la capacité du modèle à discriminer les individus en bonne et mauvaise santé. Sous SAS, ce modèle est estimé via \texttt{PROC\ LOGISTIC} avec les options \texttt{RSQUARE}, \texttt{LACKFIT} et \texttt{CTABLE} (cf.~Annexe\textasciitilde B).

\begin{table}[!h]
\centering
\caption{\label{tab:tab-logistic-binaire}Régression logistique binaire --- Coefficients et Odds Ratios.}
\centering
\resizebox{\ifdim\width>\linewidth\linewidth\else\width\fi}{!}{
\fontsize{7}{9}\selectfont
\begin{tabular}[t]{llrrrrrl}
\toprule
  & Variable & $\hat{\beta}$ & Err. std. & $z$ & p-value & OR & \\
\midrule
(Intercept) & (Intercept) & -6.4710 & 0.9325 & -6.94 & < 0.001 & 0.002 & ***\\
age & age & 0.0285 & 0.0081 & 3.52 & < 0.001 & 1.029 & ***\\
nombre_consultations_annuelles & nombre_consultations_annuelles & 0.0784 & 0.0487 & 1.61 & 0.1073 & 1.082 & \\
nombre_medicaments_reguliers & nombre_medicaments_reguliers & 0.0846 & 0.0761 & 1.11 & 0.2658 & 1.088 & \\
sexeHomme & sexeHomme & 0.1044 & 0.2546 & 0.41 & 0.6817 & 1.110 & \\
\addlinespace
hospitalisation_12_moisOui & hospitalisation_12_moisOui & 0.4827 & 0.2576 & 1.87 & 0.0610 & 1.620 & \\
activite_physiqueActivité légère & activite_physiqueActivité légère & 2.0824 & 0.7872 & 2.65 & 0.0082 & 8.023 & **\\
activite_physiqueActivité modérée & activite_physiqueActivité modérée & 0.8804 & 0.8478 & 1.04 & 0.2991 & 2.412 & \\
activite_physiqueSédentaire & activite_physiqueSédentaire & 2.9443 & 0.7758 & 3.80 & < 0.001 & 18.998 & ***\\
handicap_declareOui & handicap_declareOui & 1.4641 & 0.2626 & 5.58 & < 0.001 & 4.324 & ***\\
\bottomrule
\end{tabular}}
\end{table}

\begin{figure}[H]

{\centering \includegraphics[width=0.9\linewidth]{rapport_files/figure-latex/fig-roc-1} 

}

\caption{Courbe ROC du modèle logistique binaire. L'AUC mesure la capacité discriminante du modèle.}\label{fig:fig-roc}
\end{figure}

\textbf{Interprétation de l'AUC.} L'AUC (aire sous la courbe ROC) est de \textbf{0.864}. Selon la grille de \textcite{hosmer2013} : un AUC \(\in [0.7, 0.8]\) indique une discrimination acceptable, un AUC \(\in [0.8, 0.9]\) une bonne discrimination. Ce résultat signifie que le modèle parvient à distinguer correctement les individus en bonne et mauvaise santé dans 86~\% des cas.

\emph{Implication pour la santé publique} : un modèle prédictif performant pourrait être utilisé pour le \textbf{ciblage des populations à risque} dans les campagnes de prévention. Par exemple, les caisses d'assurance maladie pourraient identifier les profils à haut risque de mauvaise santé perçue et leur proposer des programmes de prévention personnalisés (bilans de santé, accompagnement à l'activité physique).

\emph{Implication économique} : le ciblage précis des populations à risque permettrait d'optimiser l'allocation des ressources de prévention, en concentrant les dépenses sur les individus les plus susceptibles de bénéficier d'une intervention --- un argument en faveur de la médecine personnalisée et préventive.

\hypertarget{matrice-de-confusion}{%
\subsubsection{Matrice de confusion}\label{matrice-de-confusion}}

\begin{table}[!h]
\centering
\caption{\label{tab:tab-confusion}Matrice de confusion --- Modèle logistique binaire, seuil = 0.5.}
\centering
\fontsize{9}{11}\selectfont
\begin{tabular}[t]{lrr}
\toprule
  & 0 & 1\\
\midrule
0 & 338 & 34\\
1 & 62 & 66\\
\bottomrule
\end{tabular}
\end{table}

\begin{table}[!h]
\centering
\caption{\label{tab:tab-metriques}Métriques de classification du modèle logistique binaire.}
\centering
\fontsize{9}{11}\selectfont
\begin{tabular}[t]{lr}
\toprule
Métrique & Valeur\\
\midrule
Exactitude (Accuracy) & 0.808\\
Sensibilité & 0.516\\
Spécificité & 0.909\\
\bottomrule
\end{tabular}
\end{table}

\textbf{Interprétation et implications.} La matrice de confusion (table\textasciitilde\ref{tab:tab-confusion}) et les métriques associées (table\textasciitilde\ref{tab:tab-metriques}) permettent d'évaluer la qualité de la classification. En santé publique, la \textbf{sensibilité} (capacité à détecter correctement les individus en mauvaise santé) est particulièrement importante, car un faux négatif signifie qu'un individu à risque n'est pas identifié. La \textbf{spécificité} (capacité à identifier correctement les individus en bonne santé) est importante pour éviter les interventions inutiles (faux positifs) et les coûts associés.

\emph{Implication sociétale} : un outil de détection précoce de la mauvaise santé perçue, fondé sur des variables facilement mesurables (âge, consultations, médicaments, activité physique), pourrait être intégré dans les systèmes d'information des caisses d'assurance maladie pour alimenter les politiques de prévention ciblée.

\bigskip

\textbf{Conclusion du chapitre.} L'ensemble des résultats --- analyse bivariée, régression linéaire, ANCOVA, régression logistique ordinale et binaire --- converge vers un \textbf{noyau dur de déterminants robustes} de la santé perçue : l'âge, les consultations médicales, les médicaments réguliers, l'activité physique et le handicap déclaré. Ces résultats sont cohérents avec le modèle de capital santé de Grossman, le gradient social de Marmot, et les travaux empiriques sur la santé perçue en France. Chaque résultat a été interprété au regard de ses implications économiques, sanitaires, sociodémographiques et sociétales.

\hypertarget{chap-robustesse}{%
\section{Robustesse, limites et discussion}\label{chap-robustesse}}

Ce chapitre discute la robustesse des résultats obtenus, leur conformité à la réalité épidémiologique, et identifie les limites du travail. Chaque point est assorti de ses implications pour la recherche et les politiques publiques.

\hypertarget{synthuxe8se-des-ruxe9sultats-et-conformituxe9-uxe0-la-ruxe9alituxe9}{%
\subsection{Synthèse des résultats et conformité à la réalité}\label{synthuxe8se-des-ruxe9sultats-et-conformituxe9-uxe0-la-ruxe9alituxe9}}

Les analyses convergent vers un ensemble de déterminants robustes de l'état de santé perçu. L'âge, le nombre de consultations médicales, le nombre de médicaments réguliers, l'activité physique et le handicap déclaré sont systématiquement significatifs, quelle que soit la méthode employée (régression linéaire, ANCOVA, régression logistique ordinale et binaire).

\textbf{Ces résultats sont-ils conformes à la réalité ?} La réponse est positive sur trois plans :

\begin{enumerate}
\def\labelenumi{\arabic{enumi}.}
\tightlist
\item
  \textbf{Conformité épidémiologique} : la dégradation de la santé perçue avec l'âge est un résultat universel \autocite{mackenbach1999}. L'effet du handicap est conforme aux travaux de \textcite{cambois2008} sur la « double peine ».
\item
  \textbf{Conformité au modèle théorique} : l'association consultations/médicaments--mauvaise santé perçue est cohérente avec le modèle de Grossman \autocite{grossman1972}.
\item
  \textbf{Conformité méthodologique} : la convergence des méthodes de sélection (Forward, Backward, Stepwise) et la concordance entre les modèles linéaire et ordinal témoignent de la stabilité.
\end{enumerate}

\emph{Implication pour les politiques publiques} : la convergence des résultats renforce la crédibilité des recommandations. Les décideurs peuvent s'appuyer sur des résultats robustes à la spécification pour concevoir des politiques de prévention ciblées (seniors, sédentaires, personnes handicapées, polymédiqués).

\hypertarget{discussion-des-hypothuxe8ses}{%
\subsection{Discussion des hypothèses}\label{discussion-des-hypothuxe8ses}}

\hypertarget{hypothuxe8se-h1-duxe9terminants-quantitatifs-dominants}{%
\subsubsection{Hypothèse H1 --- Déterminants quantitatifs dominants}\label{hypothuxe8se-h1-duxe9terminants-quantitatifs-dominants}}

L'âge, le nombre de consultations et le nombre de médicaments sont les trois variables quantitatives systématiquement retenues par toutes les procédures de sélection. Le VIF ne révèle pas de colinéarité problématique. \textbf{H1 est validée.}

\emph{Implication économique} : le vieillissement démographique (papy-boom, allongement de l'espérance de vie) va mécaniquement augmenter la proportion d'individus se déclarant en mauvaise santé, avec des conséquences sur les dépenses de l'assurance maladie et la soutenabilité du système de retraite.

\hypertarget{hypothuxe8se-h2-effet-propre-des-facteurs-comportementaux}{%
\subsubsection{Hypothèse H2 --- Effet propre des facteurs comportementaux}\label{hypothuxe8se-h2-effet-propre-des-facteurs-comportementaux}}

L'ANCOVA montre que l'activité physique et le handicap déclaré ont un effet significatif sur le score de santé, au-delà des effets sociodémographiques. \textbf{H2 est validée.}

\emph{Implication sociétale} : l'activité physique est un levier d'action modifiable, contrairement à l'âge ou au sexe. La politique de prévention par l'activité physique (sport-santé) est l'une des rares interventions dont l'effet sur la santé perçue est démontré même après ajustement sur les facteurs confondants.

\hypertarget{hypothuxe8se-h3-supuxe9riorituxe9-du-moduxe8le-ordinal}{%
\subsubsection{Hypothèse H3 --- Supériorité du modèle ordinal}\label{hypothuxe8se-h3-supuxe9riorituxe9-du-moduxe8le-ordinal}}

La comparaison entre le modèle linéaire et le modèle logistique ordinal montre que les deux approches identifient les mêmes déterminants significatifs, confortant la robustesse des conclusions.

\hypertarget{limites-de-donnuxe9es}{%
\subsection{Limites de données}\label{limites-de-donnuxe9es}}

\begin{itemize}
\tightlist
\item
  \textbf{Taille de l'échantillon} : \(n = 500\) est suffisant pour les analyses bivariées et la régression, mais limite la puissance des tests de sous-groupes.
\item
  \textbf{Variables omises} : des déterminants importants (alimentation, qualité du sommeil, pollution, capital social, santé mentale) ne sont pas disponibles.
\end{itemize}

\hypertarget{limites-muxe9thodologiques}{%
\subsection{Limites méthodologiques}\label{limites-muxe9thodologiques}}

\begin{itemize}
\tightlist
\item
  \textbf{Causalité} : le design en coupe transversale interdit toute interprétation causale.
\item
  \textbf{Endogénéité} : le nombre de consultations et le nombre de médicaments sont potentiellement endogènes \autocite{grossman1972}.
\item
  \textbf{Hypothèse d'équidistance} : la conversion en score ordinal (1--5) pour la régression linéaire suppose une équidistance non vérifiable.
\end{itemize}

\hypertarget{ouvertures}{%
\subsection{Ouvertures}\label{ouvertures}}

\begin{enumerate}
\def\labelenumi{\arabic{enumi}.}
\tightlist
\item
  \textbf{Extension à d'autres vagues ESPS} : application de la même méthodologie aux données de vagues antérieures pour une analyse longitudinale.
\item
  \textbf{Modèles multiniveaux} : prise en compte de la structure régionale via \texttt{PROC\ MIXED} ou \texttt{PROC\ GLIMMIX}.
\item
  \textbf{Variables instrumentales} : identification d'instruments pour traiter l'endogénéité.
\item
  \textbf{Analyses longitudinales} : exploitation de la dimension panel de l'ESPS pour des modèles à effets fixes.
\end{enumerate}

\bigskip

\textbf{Conclusion du chapitre.} Les résultats sont robustes à la spécification et cohérents avec la littérature. Les limites identifiées appellent à la prudence mais n'invalident pas l'exercice méthodologique.

\newpage

\hypertarget{conclusion-guxe9nuxe9rale}{%
\section*{Conclusion générale}\label{conclusion-guxe9nuxe9rale}}
\addcontentsline{toc}{section}{Conclusion générale}

\hypertarget{bilan}{%
\subsection*{Bilan}\label{bilan}}
\addcontentsline{toc}{subsection}{Bilan}

Ce projet a mis en œuvre une analyse statistique complète et systématique des déterminants de l'état de santé perçu, à partir des données de l'Enquête Santé et Protection Sociale (ESPS). L'ensemble des traitements statistiques a été réalisé à l'aide du logiciel SAS.

\textbf{Contribution méthodologique.} La valeur ajoutée de ce travail réside dans la confrontation systématique des approches linéaire et ordinale sur un même jeu de données, assortie d'un appareil diagnostique complet (VIF, Shapiro-Wilk, Durbin-Watson, test de White, Cook, LOESS, test du Score, ROC/AUC) rarement déployé conjointement dans les études existantes. Chaque test a été justifié par sa fonction dans la chaîne de validation du modèle, et chaque résultat a été interprété au regard de ses implications économiques, sanitaires et sociétales.

\textbf{Résultats substantiels.} L'âge, le nombre de consultations médicales, le nombre de médicaments réguliers, l'activité physique et le handicap déclaré émergent comme déterminants robustes de l'état de santé perçu, significatifs à travers les différentes spécifications. L'ANCOVA confirme l'existence d'effets propres des facteurs comportementaux et de santé au-delà des caractéristiques sociodémographiques. Ces résultats sont conformes à la littérature épidémiologique \autocites[ ]{mackenbach1999}[ ]{marmot2005}{singh2006} et aux prédictions du modèle de capital santé de \textcite{grossman1972}.

\textbf{Implications pour les politiques publiques.} Les résultats de ce travail soutiennent trois axes d'action :

\begin{enumerate}
\def\labelenumi{\arabic{enumi}.}
\tightlist
\item
  \textbf{Prévention du vieillissement en mauvaise santé} : le gradient âge--santé perçue justifie les politiques de prévention ciblées sur les seniors (plans Bien vieillir, consultations de prévention).
\item
  \textbf{Promotion de l'activité physique} : l'effet propre de l'activité physique, démontré par l'ANCOVA, légitime les politiques sport-santé et les prescriptions d'activité physique adaptée.
\item
  \textbf{Accompagnement du handicap} : la « double peine » des personnes handicapées appelle des dispositifs renforcés d'accessibilité aux soins et d'accompagnement social.
\end{enumerate}

\hypertarget{perspectives}{%
\subsection*{Perspectives}\label{perspectives}}
\addcontentsline{toc}{subsection}{Perspectives}

Le design en coupe transversale et la taille de l'échantillon (\(n = 500\)) appellent à la prudence. Plusieurs pistes d'approfondissement se dessinent :

\begin{enumerate}
\def\labelenumi{\arabic{enumi}.}
\tightlist
\item
  \textbf{Extension à d'autres vagues ESPS} : l'application de la même méthodologie aux données de vagues antérieures permettrait une analyse longitudinale des tendances.
\item
  \textbf{Modèles multiniveaux} : la prise en compte de la structure régionale via \texttt{PROC\ MIXED} ou \texttt{PROC\ GLIMMIX} enrichirait l'analyse.
\item
  \textbf{Variables instrumentales} : l'identification d'instruments pour traiter l'endogénéité des variables de recours aux soins renforcerait la portée causale des résultats.
\item
  \textbf{Analyses longitudinales} : l'exploitation de la dimension panel de l'ESPS autoriserait des modèles à effets fixes.
\end{enumerate}

En définitive, ce travail démontre que la régression linéaire, convenablement diagnostiquée et confrontée au modèle ordinal, constitue un outil pertinent et interprétable pour la modélisation de la santé perçue --- un résultat méthodologique transposable à d'autres enquêtes de population et à d'autres contextes institutionnels.

\newpage

(APPENDIX) Annexes \{-\}

\hypertarget{annexe-a-description-duxe9tailluxe9e-de-la-base-de-donnuxe9es}{%
\section*{Annexe A : Description détaillée de la base de données}\label{annexe-a-description-duxe9tailluxe9e-de-la-base-de-donnuxe9es}}
\addcontentsline{toc}{section}{Annexe A : Description détaillée de la base de données}

La base \texttt{base\_esps.csv} comporte 500 observations et 25 variables. Le dictionnaire complet des variables est présenté ci-dessous.

\hypertarget{ruxe9fuxe9rence-muxe9thodologique}{%
\subsection*{Référence méthodologique}\label{ruxe9fuxe9rence-muxe9thodologique}}
\addcontentsline{toc}{subsection}{Référence méthodologique}

\begin{tabular}{@{} l p{10cm} @{}}
\toprule
\textbf{Élément} & \textbf{Détail} \\
\midrule
Source & IRDES --- Enquête Santé et Protection Sociale (ESPS) \\
Méthodologie & Enquête en coupe transversale (\textit{cross-sectional survey}) \\
Unité statistique & Individu résidant en France métropolitaine \\
Fichier & \texttt{base\_esps.csv} \\
Encodage & UTF-8 \\
Séparateur & Point-virgule (\texttt{;}) \\
Observations & 500 individus \\
Variables & 25 \\
\bottomrule
\end{tabular}

\medskip

\hypertarget{dictionnaire-des-variables}{%
\subsection*{Dictionnaire des variables}\label{dictionnaire-des-variables}}
\addcontentsline{toc}{subsection}{Dictionnaire des variables}

\begingroup
\scriptsize
\begin{longtable}{@{} c l l p{3.8cm} p{5.5cm} @{}}
\toprule
\textbf{\#} & \textbf{Variable} & \textbf{Type} & \textbf{Description} & \textbf{Modalités / Unité} \\
\midrule
\endfirsthead
\toprule
\textbf{\#} & \textbf{Variable} & \textbf{Type} & \textbf{Description} & \textbf{Modalités / Unité} \\
\midrule
\endhead
\bottomrule
\endfoot
1  & \texttt{id}                             & Identifiant  & Identifiant unique de l'individu                     & Entier (1 $\rightarrow$ 500) \\
2  & \texttt{age}                            & Quantitative & Âge de l'individu                                    & 18 -- 95 ans \\
3  & \texttt{sexe}                           & Qualitative  & Sexe                                                 & Homme, Femme \\
4  & \texttt{profession}                     & Qualitative  & Catégorie socioprofessionnelle                       & Agriculteur, Artisan/Commerçant, Cadre supérieur, Profession intermédiaire, Employé, Ouvrier, Retraité, Sans activité professionnelle, Étudiant \\
5  & \texttt{niveau\_education}              & Qualitative  & Niveau de diplôme le plus élevé                      & Sans diplôme, CAP/BEP, Baccalauréat, Bac+2, Bac+3/Licence, Bac+5/Master, Doctorat \\
6  & \texttt{situation\_familiale}           & Qualitative  & Situation matrimoniale                               & Célibataire, Marié(e), Pacsé(e), Divorcé(e), Veuf/Veuve \\
7  & \texttt{nombre\_enfants}                & Quantitative & Nombre d'enfants à charge                            & 0 -- 5 \\
8  & \texttt{region}                         & Qualitative  & Région de résidence                                  & 13 régions métropolitaines \\
9  & \texttt{type\_commune}                  & Qualitative  & Taille de la commune de résidence                    & Rurale, Urbaine petite, Urbaine moyenne, Métropole \\
10 & \texttt{revenus\_mensuels\_euros}       & Quantitative & Revenus mensuels nets (€)                            & Distribution gaussienne centrée selon la profession \\
11 & \texttt{statut\_assurance}              & Qualitative  & Régime d'assurance maladie obligatoire               & Régime général, MSA, RSI/SSI, Régime spécial, Autre régime \\
12 & \texttt{couverture\_complementaire}     & Qualitative  & Type de couverture complémentaire santé              & Mutuelle privée, CSS, Assurance privée, Aucune complémentaire \\
13 & \texttt{depenses\_sante\_annuelles\_euros} & Quantitative & Dépenses de santé annuelles (€)                   & Corrélées positivement à l'âge \\
14 & \texttt{nombre\_consultations\_annuelles} & Quantitative & Nombre de consultations médicales sur 12 mois      & 0 -- 30 \\
15 & \texttt{nombre\_medicaments\_reguliers} & Quantitative & Nombre de médicaments pris régulièrement             & Corrélé à l'âge \\
16 & \texttt{hospitalisation\_12\_mois}      & Qualitative  & Hospitalisation au cours des 12 derniers mois        & Oui, Non \\
17 & \texttt{etat\_sante\_percu}             & Qualitative  & Auto-évaluation de l'état de santé                   & Très bon, Bon, Assez bon, Mauvais, Très mauvais \\
18 & \texttt{maladies\_declarees}            & Qualitative  & Maladie chronique principale déclarée                & Aucune, Diabète, Hypertension, Asthme, Dépression, Lombalgie chronique, Arthrose, Cardiopathie, Cancer (rémission), Maladie respiratoire chronique, Trouble anxieux, Migraine chronique, Obésité morbide \\
19 & \texttt{imc}                            & Quantitative & Indice de masse corporelle                           & 15.0 -- 50.0 (moyenne $\approx$ 25.5) \\
20 & \texttt{tabagisme}                      & Qualitative  & Statut tabagique                                     & Non-fumeur, Ancien fumeur, Fumeur occasionnel, Fumeur quotidien \\
21 & \texttt{consommation\_alcool}           & Qualitative  & Fréquence de consommation d'alcool                   & Jamais, Occasionnelle, Régulière, Quotidienne \\
22 & \texttt{activite\_physique}             & Qualitative  & Niveau d'activité physique                           & Sédentaire, Activité légère, Activité modérée, Activité intense \\
23 & \texttt{handicap\_declare}              & Qualitative  & Handicap ou limitation déclaré(e)                    & Oui, Non \\
24 & \texttt{renoncement\_soins\_12\_mois}   & Qualitative  & Renoncement à des soins pour raisons financières     & Oui, Non \\
25 & \texttt{satisfaction\_systeme\_soins}    & Qualitative  & Satisfaction vis-à-vis du système de soins           & Très satisfait, Satisfait, Peu satisfait, Pas du tout satisfait \\
\end{longtable}
\endgroup

\hypertarget{ruxe9partition-des-types-de-variables}{%
\subsection*{Répartition des types de variables}\label{ruxe9partition-des-types-de-variables}}
\addcontentsline{toc}{subsection}{Répartition des types de variables}

\begin{tabular}{@{} l c p{11cm} @{}}
\toprule
\textbf{Type} & \textbf{Nombre} & \textbf{Variables} \\
\midrule
Identifiant   & 1  & \texttt{id} \\
Quantitatives & 7  & \texttt{age}, \texttt{nombre\_enfants}, \texttt{revenus\_mensuels\_euros}, \texttt{depenses\_sante\_annuelles\_euros}, \texttt{nombre\_consultations\_annuelles}, \texttt{nombre\_medicaments\_reguliers}, \texttt{imc} \\
Qualitatives  & 17 & \texttt{sexe}, \texttt{profession}, \texttt{niveau\_education}, \texttt{situation\_familiale}, \texttt{region}, \texttt{type\_commune}, \texttt{statut\_assurance}, \texttt{couverture\_complementaire}, \texttt{hospitalisation\_12\_mois}, \texttt{etat\_sante\_percu}, \texttt{maladies\_declarees}, \texttt{tabagisme}, \texttt{consommation\_alcool}, \texttt{activite\_physique}, \texttt{handicap\_declare}, \texttt{renoncement\_soins\_12\_mois}, \texttt{satisfaction\_systeme\_soins} \\
\bottomrule
\end{tabular}

\hypertarget{corruxe9lations-intuxe9gruxe9es}{%
\subsection*{Corrélations intégrées}\label{corruxe9lations-intuxe9gruxe9es}}
\addcontentsline{toc}{subsection}{Corrélations intégrées}

Les principales corrélations observées dans les données sont conformes aux associations documentées dans les enquêtes de santé françaises (ESPS/DREES) :

\begin{itemize}
\tightlist
\item
  \textbf{Profession \(\leftrightarrow\) Âge} : les retraités sont majoritairement \(\geq\) 65 ans ; les étudiants \(\leq\) 25 ans.
\item
  \textbf{Revenus \(\leftrightarrow\) Profession} : les cadres ont des revenus plus élevés que les ouvriers.
\item
  \textbf{Dépenses de santé \(\leftrightarrow\) Âge} : les dépenses augmentent avec l'âge.
\item
  \textbf{Consultations \(\leftrightarrow\) Âge} : le nombre de consultations croît avec l'âge.
\item
  \textbf{État de santé perçu \(\leftrightarrow\) Âge} : dégradation perçue avec l'avancée en âge.
\item
  \textbf{Handicap \(\leftrightarrow\) Âge} : probabilité de handicap déclaré croissante avec l'âge.
\item
  \textbf{Renoncement aux soins \(\leftrightarrow\) Revenus} : le renoncement est plus fréquent chez les bas revenus.
\end{itemize}

\hypertarget{annexe-b-code-sas-complet}{%
\section*{Annexe B : Code SAS complet}\label{annexe-b-code-sas-complet}}
\addcontentsline{toc}{section}{Annexe B : Code SAS complet}

L'intégralité du code SAS utilisé pour les traitements statistiques est reproduite ci-dessous.

\begin{footnotesize}
\begin{verbatim}
/*****************************************************************************
 * Analyses statistiques complètes — Base ESPS
 * Variable réponse : etat_sante_percu (qualitative ordinale)
 * Modalités : Très bon | Bon | Assez bon | Mauvais | Très mauvais
 *
 * Contexte : Enquête Santé et Protection Sociale (ESPS)
 * Fichier : base_esps.csv (500 obs. × 25 variables)
 *
 * Analyses couvertes :
 * 0. Importation et recodage
 * 1. Analyse bivariée
 * 2. Régression linéaire et R²
 * 3. Étude de la colinéarité
 * 4. Sélection de variables
 * 5. Validation du modèle
 * 6. Analyse de covariance (ANCOVA)
 * 7. Régression logistique ordinale
 *****************************************************************************/


/* ═══════════════════════════════════════════════════════════════════════════
 0. IMPORTATION ET RECODAGE
 ═══════════════════════════════════════════════════════════════════════════ */

/* Définir la bibliothèque de travail */
libname proj "/home/u64471098/MonProjet";

proc import datafile="/home/u64471098/MonProjet/base_esps.csv"
 out = proj.base_esps
 dbms = dlm
 replace;
 delimiter = ";";
 getnames = yes;
 guessingrows = 500;
run;

/* Vérification rapide */
proc contents data = proj.base_esps;
run;

proc freq data = proj.base_esps;
 tables etat_sante_percu;
 title "Distribution de la variable réponse — État de santé perçu";
run;

/* -----------------------------------------------------------------------
 Recodage numérique ordinal de etat_sante_percu
 1 = Très mauvais → 5 = Très bon (sens croissant = meilleure santé)
 ----------------------------------------------------------------------- */
data proj.base_esps;
 set proj.base_esps;

 /* Score numérique ordinal */
 if etat_sante_percu = "Très mauvais" then score_sante = 1;
 else if etat_sante_percu = "Mauvais" then score_sante = 2;
 else if etat_sante_percu = "Assez bon" then score_sante = 3;
 else if etat_sante_percu = "Bon" then score_sante = 4;
 else if etat_sante_percu = "Très bon" then score_sante = 5;
 else score_sante =.; /* valeur manquante si modalité inattendue */

 /* Variable binaire pour la régression logistique binaire (optionnel) :
 1 = Mauvaise santé (Très mauvais + Mauvais)
 0 = Bonne santé (Assez bon + Bon + Très bon) */
 if score_sante in (1, 2) then sante_mauvaise = 1;
 else sante_mauvaise = 0;

 /* Format pour hospitalisation et handicap */
 if hospitalisation_12_mois = "Oui" then hospit_num = 1;
 else hospit_num = 0;
 if handicap_declare = "Oui" then handicap_num = 1;
 else handicap_num = 0;
 if renoncement_soins_12_mois = "Oui" then renoncement_num = 1;
 else renoncement_num = 0;
run;

/* Vérifier le recodage */
proc freq data = proj.base_esps;
 tables etat_sante_percu * score_sante / norow nocol nopercent;
 title "Vérification du recodage ordinal";
run;


/* ═══════════════════════════════════════════════════════════════════════════
 1. ANALYSE BIVARIÉE
 ═══════════════════════════════════════════════════════════════════════════
 Y = etat_sante_percu (qualitative ordinale)
 → Y quali × X quali : test du Chi-2 (PROC FREQ)
 → Y quali × X quanti : ANOVA ou Kruskal-Wallis (PROC NPAR1WAY)
 ═══════════════════════════════════════════════════════════════════════════ */

/* 1a. Croisement avec les variables qualitatives — Test du Chi-2 */

proc freq data = proj.base_esps;
 tables etat_sante_percu * sexe / chisq expected cellchi2 norow nocol;
 title "Bivariée : État de santé perçu × Sexe";
run;

proc freq data = proj.base_esps;
 tables etat_sante_percu * profession / chisq expected;
 title "Bivariée : État de santé perçu × Profession";
run;

proc freq data = proj.base_esps;
 tables etat_sante_percu * niveau_education / chisq expected;
 title "Bivariée : État de santé perçu × Niveau d'éducation";
run;

proc freq data = proj.base_esps;
 tables etat_sante_percu * situation_familiale / chisq expected;
 title "Bivariée : État de santé perçu × Situation familiale";
run;

proc freq data = proj.base_esps;
 tables etat_sante_percu * couverture_complementaire / chisq expected;
 title "Bivariée : État de santé perçu × Couverture complémentaire";
run;

proc freq data = proj.base_esps;
 tables etat_sante_percu * hospitalisation_12_mois / chisq expected;
 title "Bivariée : État de santé perçu × Hospitalisation 12 mois";
run;

proc freq data = proj.base_esps;
 tables etat_sante_percu * tabagisme / chisq expected;
 title "Bivariée : État de santé perçu × Tabagisme";
run;

proc freq data = proj.base_esps;
 tables etat_sante_percu * consommation_alcool / chisq expected;
 title "Bivariée : État de santé perçu × Consommation d'alcool";
run;

proc freq data = proj.base_esps;
 tables etat_sante_percu * activite_physique / chisq expected;
 title "Bivariée : État de santé perçu × Activité physique";
run;

proc freq data = proj.base_esps;
 tables etat_sante_percu * handicap_declare / chisq expected;
 title "Bivariée : État de santé perçu × Handicap déclaré";
run;

proc freq data = proj.base_esps;
 tables etat_sante_percu * renoncement_soins_12_mois / chisq expected;
 title "Bivariée : État de santé perçu × Renoncement aux soins";
run;

proc freq data = proj.base_esps;
 tables etat_sante_percu * maladies_declarees / chisq expected;
 title "Bivariée : État de santé perçu × Maladies déclarées";
run;

/* 1b. Croisement avec les variables quantitatives
 ANOVA à un facteur : Y = etat_sante_percu (facteur),
 X = variable quanti */

proc glm data = proj.base_esps;
 class etat_sante_percu;
 model age = etat_sante_percu;
 means etat_sante_percu / tukey;
 title "Bivariée ANOVA : Âge selon l'état de santé perçu";
run; quit;

proc glm data = proj.base_esps;
 class etat_sante_percu;
 model revenus_mensuels_euros = etat_sante_percu;
 means etat_sante_percu / tukey;
 title "Bivariée ANOVA : Revenus mensuels selon l'état de santé perçu";
run; quit;

proc glm data = proj.base_esps;
 class etat_sante_percu;
 model imc = etat_sante_percu;
 means etat_sante_percu / tukey;
 title "Bivariée ANOVA : IMC selon l'état de santé perçu";
run; quit;

proc glm data = proj.base_esps;
 class etat_sante_percu;
 model depenses_sante_annuelles_euros = etat_sante_percu;
 means etat_sante_percu / tukey;
 title "Bivariée ANOVA : Dépenses de santé selon l'état de santé perçu";
run; quit;

proc glm data = proj.base_esps;
 class etat_sante_percu;
 model nombre_consultations_annuelles = etat_sante_percu;
 means etat_sante_percu / tukey;
 title "Bivariée ANOVA : Consultations selon l'état de santé perçu";
run; quit;

proc glm data = proj.base_esps;
 class etat_sante_percu;
 model nombre_medicaments_reguliers = etat_sante_percu;
 means etat_sante_percu / tukey;
 title "Bivariée ANOVA : Médicaments réguliers selon l'état de santé perçu";
run; quit;

/* 1c. Alternative non paramétrique — Test de Kruskal-Wallis */

proc npar1way data = proj.base_esps wilcoxon;
 class etat_sante_percu;
 var age revenus_mensuels_euros imc depenses_sante_annuelles_euros
 nombre_consultations_annuelles nombre_medicaments_reguliers;
 title "Test de Kruskal-Wallis : Variables quantitatives
 selon l'état de santé perçu";
run;

/* 1d. Corrélation de Spearman */
proc corr data = proj.base_esps spearman;
 var score_sante;
 with age revenus_mensuels_euros imc depenses_sante_annuelles_euros
 nombre_consultations_annuelles nombre_medicaments_reguliers
 nombre_enfants;
 title "Corrélation de Spearman : score_sante vs variables quantitatives";
run;

/* 1e. Graphiques bivariés */
proc sgplot data = proj.base_esps;
 vbox age / category = etat_sante_percu;
 title "Boîte à moustaches : Âge selon l'état de santé perçu";
run;

proc sgplot data = proj.base_esps;
 vbox imc / category = etat_sante_percu;
 title "Boîte à moustaches : IMC selon l'état de santé perçu";
run;

proc sgplot data = proj.base_esps;
 vbox depenses_sante_annuelles_euros / category = etat_sante_percu;
 title "Boîte à moustaches : Dépenses de santé selon l'état de santé perçu";
run;

proc sgplot data = proj.base_esps;
 vbox nombre_consultations_annuelles / category = etat_sante_percu;
 title "Boîte à moustaches : Consultations selon l'état de santé perçu";
run;


/* ═══════════════════════════════════════════════════════════════════════════
 2. RÉGRESSION LINÉAIRE ET R²
 ═══════════════════════════════════════════════════════════════════════════ */

/* 2a. Modèle complet — variables quantitatives */
proc reg data = proj.base_esps;
 model score_sante =
 age
 revenus_mensuels_euros
 depenses_sante_annuelles_euros
 nombre_consultations_annuelles
 nombre_medicaments_reguliers
 imc
 nombre_enfants;
 title "Régression linéaire — Modèle complet (score_sante)";
run; quit;

/* 2b. Modèle réduit — variables significatives uniquement */
proc reg data = proj.base_esps;
 model score_sante =
 age
 nombre_consultations_annuelles
 nombre_medicaments_reguliers;
 title "Régression linéaire — Modèle réduit (score_sante)";
run; quit;


/* ═══════════════════════════════════════════════════════════════════════════
 3. ÉTUDE DE LA COLINÉARITÉ
 ═══════════════════════════════════════════════════════════════════════════ */

/* 3a. VIF sur toutes les variables quantitatives */
proc reg data = proj.base_esps;
 model score_sante =
 age
 revenus_mensuels_euros
 depenses_sante_annuelles_euros
 nombre_consultations_annuelles
 nombre_medicaments_reguliers
 imc
 nombre_enfants / vif tol;
 title "Colinéarité — VIF et Tolérance (score_sante)";
run; quit;

/* 3b. Matrice de corrélation entre les prédicteurs */
proc corr data = proj.base_esps;
 var age
 revenus_mensuels_euros
 depenses_sante_annuelles_euros
 nombre_consultations_annuelles
 nombre_medicaments_reguliers
 imc
 nombre_enfants;
 title "Matrice de corrélation — Variables explicatives quantitatives";
run;

/* 3c. Indices de conditionnement (collin) */
proc reg data = proj.base_esps;
 model score_sante =
 age
 revenus_mensuels_euros
 depenses_sante_annuelles_euros
 nombre_consultations_annuelles
 nombre_medicaments_reguliers
 imc
 nombre_enfants / collin;
 title "Indices de conditionnement (collinéarité diagnostics)";
run; quit;


/* ═══════════════════════════════════════════════════════════════════════════
 4. SÉLECTION DE VARIABLES
 ═══════════════════════════════════════════════════════════════════════════ */

/* 4a. Sélection par R² ajusté, AIC et BIC */
proc reg data = proj.base_esps;
 model score_sante =
 age
 revenus_mensuels_euros
 depenses_sante_annuelles_euros
 nombre_consultations_annuelles
 nombre_medicaments_reguliers
 imc
 nombre_enfants / selection = ADJRSQ aic bic best = 5;
 title "Sélection de modèles — R² ajusté, AIC, BIC";
run; quit;

/* 4b. Sélection ascendante (forward) */
proc reg data = proj.base_esps;
 model score_sante =
 age
 revenus_mensuels_euros
 depenses_sante_annuelles_euros
 nombre_consultations_annuelles
 nombre_medicaments_reguliers
 imc
 nombre_enfants / selection = forward slentry = 0.05;
 title "Sélection Forward (seuil d'entrée = 0.05)";
run; quit;

/* 4c. Sélection descendante (backward) */
proc reg data = proj.base_esps;
 model score_sante =
 age
 revenus_mensuels_euros
 depenses_sante_annuelles_euros
 nombre_consultations_annuelles
 nombre_medicaments_reguliers
 imc
 nombre_enfants / selection = backward slstay = 0.05;
 title "Sélection Backward (seuil de maintien = 0.05)";
run; quit;

/* 4d. Sélection pas à pas (stepwise) */
proc reg data = proj.base_esps;
 model score_sante =
 age
 revenus_mensuels_euros
 depenses_sante_annuelles_euros
 nombre_consultations_annuelles
 nombre_medicaments_reguliers
 imc
 nombre_enfants / selection = stepwise slentry = 0.05 slstay = 0.05;
 title "Sélection Stepwise (entrée/maintien = 0.05)";
run; quit;


/* ═══════════════════════════════════════════════════════════════════════════
 5. VALIDATION DU MODÈLE
 ═══════════════════════════════════════════════════════════════════════════ */

/* 5a. Estimation du modèle retenu + export des résidus */
proc reg data = proj.base_esps
 plots(only label) = (RStudentByLeverage CooksD
 residuals(smooth) partialplot);
 model score_sante =
 age
 nombre_consultations_annuelles
 nombre_medicaments_reguliers / r spec dwprob;
 output out = proj.base_esps_resid
 rstudent = rstud
 predicted = yhat
 residual = resid
 cookd = cooksd
 h = leverage;
 title "Modèle retenu — Diagnostics complets (score_sante)";
run; quit;

/* 5b. Normalité des résidus — QQ-plot */
proc univariate data = proj.base_esps_resid normal;
 var resid;
 qqplot resid / normal(mu = 0 sigma = est);
 title "Normalité des résidus — QQ-plot";
run;

/* 5c. Normalité des résidus — Histogramme */
proc univariate data = proj.base_esps_resid;
 var resid;
 histogram resid / normal;
 title "Normalité des résidus — Histogramme";
run;

/* 5d. Éléments aberrants — RStudent vs Leverage */
proc reg data = proj.base_esps
 plots(only label) = (RStudentByLeverage);
 model score_sante =
 age
 nombre_consultations_annuelles
 nombre_medicaments_reguliers;
 title "Éléments aberrants — RStudent × Leverage";
run; quit;

/* 5e. Distance de Cook (seuil : D_i > 4/n = 4/500 = 0.008) */
proc reg data = proj.base_esps
 plots(only label) = (CooksD);
 model score_sante =
 age
 nombre_consultations_annuelles
 nombre_medicaments_reguliers;
 title "Distance de Cook — Observations influentes";
run; quit;

/* 5f. Indépendance — Test de Durbin-Watson */
proc reg data = proj.base_esps;
 model score_sante =
 age
 nombre_consultations_annuelles
 nombre_medicaments_reguliers / dwprob;
 title "Indépendance des résidus — Durbin-Watson";
run; quit;

/* 5g. Homoscédasticité — Test de White (/SPEC) */
proc reg data = proj.base_esps;
 model score_sante =
 age
 nombre_consultations_annuelles
 nombre_medicaments_reguliers / spec;
 title "Homoscédasticité — Test de White (SPEC)";
run; quit;

/* 5h. Linéarité — Résidus partiels + LOESS */
proc reg data = proj.base_esps
 plots(only) = (residuals(smooth) partialplot);
 model score_sante =
 age
 nombre_consultations_annuelles
 nombre_medicaments_reguliers / spec;
 title "Linéarité — Résidus partiels";
run; quit;

/* 5i. Visualisation LOESS */
proc sgplot data = proj.base_esps;
 loess y = score_sante x = age;
 title "LOESS — Score santé perçue vs Âge";
run;

proc sgplot data = proj.base_esps;
 loess y = score_sante x = nombre_consultations_annuelles;
 title "LOESS — Score santé perçue vs Nombre de consultations";
run;

proc sgplot data = proj.base_esps;
 loess y = score_sante x = nombre_medicaments_reguliers;
 title "LOESS — Score santé perçue vs Nombre de médicaments";
run;

/* 5j. Observations potentiellement aberrantes / influentes */
proc print data = proj.base_esps_resid;
 where abs(rstud) > 2 or cooksd > 0.008;
 var id age score_sante yhat resid rstud cooksd leverage;
 title "Observations potentiellement aberrantes / influentes";
run;


/* ═══════════════════════════════════════════════════════════════════════════
 6. ANALYSE DE COVARIANCE (ANCOVA)
 ═══════════════════════════════════════════════════════════════════════════ */

/* 6a. ANCOVA complète */
proc glm data = proj.base_esps;
 class sexe
 profession
 couverture_complementaire
 hospitalisation_12_mois
 tabagisme
 activite_physique
 handicap_declare
 maladies_declarees;
 model score_sante =
 /* Variables quantitatives (covariables) */
 age
 revenus_mensuels_euros
 imc
 depenses_sante_annuelles_euros
 nombre_consultations_annuelles
 nombre_medicaments_reguliers
 /* Variables qualitatives (facteurs) */
 sexe
 profession
 couverture_complementaire
 hospitalisation_12_mois
 tabagisme
 activite_physique
 handicap_declare
 maladies_declarees / solution ss3;
 title "ANCOVA — score_sante (modèle complet)";
run; quit;

/* 6b. ANCOVA réduite — facteurs significatifs */
proc glm data = proj.base_esps;
 class sexe
 hospitalisation_12_mois
 activite_physique
 handicap_declare;
 model score_sante =
 age
 nombre_consultations_annuelles
 nombre_medicaments_reguliers
 sexe
 hospitalisation_12_mois
 activite_physique
 handicap_declare / solution ss3;
 lsmeans sexe hospitalisation_12_mois
 activite_physique handicap_declare
 / pdiff adjust = tukey;
 title "ANCOVA réduite — score_sante (facteurs significatifs)";
run; quit;

/* 6c. Test d'interaction Âge × Sexe */
proc glm data = proj.base_esps;
 class sexe;
 model score_sante = age sexe age*sexe / solution ss3;
 title "ANCOVA — Test d'interaction Âge × Sexe sur score_sante";
run; quit;


/* ═══════════════════════════════════════════════════════════════════════════
 7. RÉGRESSION LOGISTIQUE
 ═══════════════════════════════════════════════════════════════════════════ */

/* 7a. Régression logistique ORDINALE (cumulative logit) — modèle complet */
proc logistic data = proj.base_esps descending;
 class sexe (ref = "Homme")
 hospitalisation_12_mois (ref = "Non")
 tabagisme (ref = "Non-fumeur")
 activite_physique (ref = "Sédentaire")
 handicap_declare (ref = "Non")
 couverture_complementaire (ref = "Mutuelle privée")
 / param = ref;
 model etat_sante_percu =
 age
 revenus_mensuels_euros
 imc
 nombre_consultations_annuelles
 nombre_medicaments_reguliers
 depenses_sante_annuelles_euros
 sexe
 hospitalisation_12_mois
 tabagisme
 activite_physique
 handicap_declare
 couverture_complementaire / link = cumlogit
 rsquare
 lackfit
 stb;
 title "Régression logistique ordinale — etat_sante_percu
 (modèle complet)";
run;

/* 7b. Test de l'hypothèse des cotes proportionnelles (Score Test)
 produit automatiquement par PROC LOGISTIC avec link=cumlogit.
 Si p-value < 0.05 → l'hypothèse de proportionnalité est violée. */

/* 7c. Régression logistique ordinale — modèle réduit */
proc logistic data = proj.base_esps descending;
 class sexe (ref = "Homme")
 hospitalisation_12_mois (ref = "Non")
 activite_physique (ref = "Sédentaire")
 handicap_declare (ref = "Non") / param = ref;
 model etat_sante_percu =
 age
 nombre_consultations_annuelles
 nombre_medicaments_reguliers
 sexe
 hospitalisation_12_mois
 activite_physique
 handicap_declare / link = cumlogit
 rsquare
 lackfit
 stb;
 title "Régression logistique ordinale — etat_sante_percu
 (modèle réduit)";
run;

/* 7d. Régression logistique BINAIRE
 Variable réponse : sante_mauvaise (0/1) */
proc logistic data = proj.base_esps descending;
 class sexe (ref = "Homme")
 hospitalisation_12_mois (ref = "Non")
 tabagisme (ref = "Non-fumeur")
 activite_physique (ref = "Sédentaire")
 handicap_declare (ref = "Non") / param = ref;
 model sante_mauvaise (event = "1") =
 age
 revenus_mensuels_euros
 imc
 nombre_consultations_annuelles
 nombre_medicaments_reguliers
 sexe
 hospitalisation_12_mois
 tabagisme
 activite_physique
 handicap_declare / rsquare
 lackfit
 ctable
 stb;
 title "Régression logistique binaire — sante_mauvaise (0/1)";
run;

/* 7e. Sélection de variables en régression logistique (stepwise) */
proc logistic data = proj.base_esps descending;
 class sexe (ref = "Homme")
 hospitalisation_12_mois (ref = "Non")
 tabagisme (ref = "Non-fumeur")
 activite_physique (ref = "Sédentaire")
 handicap_declare (ref = "Non")
 couverture_complementaire (ref = "Mutuelle privée")
 / param = ref;
 model etat_sante_percu =
 age
 revenus_mensuels_euros
 imc
 nombre_consultations_annuelles
 nombre_medicaments_reguliers
 depenses_sante_annuelles_euros
 sexe
 hospitalisation_12_mois
 tabagisme
 activite_physique
 handicap_declare
 couverture_complementaire / link = cumlogit
 selection = stepwise
 slentry = 0.05
 slstay = 0.05;
 title "Sélection stepwise — Régression logistique ordinale";
run;

/* 7f. Courbe ROC (régression binaire uniquement) */
proc logistic data = proj.base_esps descending
 plots(only) = (roc effect oddsratio);
 class sexe (ref = "Homme")
 hospitalisation_12_mois (ref = "Non")
 activite_physique (ref = "Sédentaire")
 handicap_declare (ref = "Non") / param = ref;
 model sante_mauvaise (event = "1") =
 age
 nombre_consultations_annuelles
 nombre_medicaments_reguliers
 sexe
 hospitalisation_12_mois
 activite_physique
 handicap_declare / rsquare
 lackfit
 ctable;
 title "Régression logistique binaire — Courbe ROC";
run;

/* 7g. Matrice de confusion (table de classement)
 → produite automatiquement par l'option ctable ci-dessus.
 Interprétation : sensibilité, spécificité, taux de classement. */
\end{verbatim}
\end{footnotesize}

\newpage

\hypertarget{annexe-c-code-r-complet}{%
\section*{Annexe C : Code R complet}\label{annexe-c-code-r-complet}}
\addcontentsline{toc}{section}{Annexe C : Code R complet}

L'intégralité du code R reproduisant fidèlement les analyses du code SAS est présentée ci-dessous.
Le script autonome est également disponible à la racine du projet sous le nom \texttt{annexe\_C\_code\_R.R}.

\medskip

\noindent\textbf{Correspondances SAS $\rightarrow$ R}

\begin{tabular}{@{} l l l @{}}
\toprule
\textbf{Section} & \textbf{Procédure SAS} & \textbf{Équivalent R} \\
\midrule
\S{}0 Importation & \texttt{PROC IMPORT} & \texttt{read.csv2()} \\
\S{}1a Bivariée quali & \texttt{PROC FREQ / chisq measures} & \texttt{chisq.test()} + \texttt{vcd::assocstats()} \\
\S{}1b ANOVA & \texttt{PROC GLM / tukey} & \texttt{aov()} + \texttt{TukeyHSD()} \\
\S{}1c Kruskal-Wallis & \texttt{PROC NPAR1WAY wilcoxon} & \texttt{kruskal.test()} \\
\S{}1d Spearman & \texttt{PROC CORR spearman} & \texttt{cor.test(method=``spearman'')} \\
\S{}2 Régr. linéaire & \texttt{PROC REG} & \texttt{lm()} \\
\S{}3 VIF & \texttt{PROC REG / vif tol} & \texttt{car::vif()} \\
\S{}3 Corrélation & \texttt{PROC CORR} & \texttt{cor()} + \texttt{corrplot::corrplot()} \\
\S{}3 Cond. number & \texttt{PROC REG / collin} & \texttt{svd()}, rapport max/min \\
\S{}4 Sous-modèles & \texttt{PROC REG / selection=ADJRSQ best=5} & \texttt{leaps::regsubsets()} \\
\S{}4 Forward & \texttt{PROC REG / selection=forward} & \texttt{step(direction=``forward'')} \\
\S{}4 Backward & \texttt{PROC REG / selection=backward} & \texttt{step(direction=``backward'')} \\
\S{}4 Stepwise & \texttt{PROC REG / selection=stepwise} & \texttt{step(direction=``both'')} \\
\S{}5 Résidus & \texttt{OUTPUT rstudent / cookd / h} & \texttt{rstudent()}, \texttt{cooks.distance()}, \texttt{hatvalues()} \\
\S{}5 Normalité & \texttt{PROC UNIVARIATE normal qqplot} & \texttt{shapiro.test()} + \texttt{qqnorm()} \\
\S{}5 Durbin-Watson & \texttt{PROC REG / dwprob} & \texttt{lmtest::dwtest()} \\
\S{}5 White & \texttt{PROC REG / spec} & \texttt{lmtest::bptest()} \\
\S{}5 Résidus partiels & \texttt{PROC REG plots(partialplot)} & \texttt{car::avPlots()} \\
\S{}5 LOESS & \texttt{PROC SGPLOT loess} & \texttt{geom\_smooth(``loess'')} \\
\S{}5k K-fold & Macro \texttt{\%kfold\_ordinal} & Boucle \texttt{MASS::polr()} + \texttt{set.seed(42)} \\
\S{}6 ANCOVA & \texttt{PROC GLM / ss3} & \texttt{lm()} + \texttt{car::Anova(type=``III'')} \\
\S{}6 LSMEANS & \texttt{LSMEANS / pdiff adjust=tukey} & \texttt{emmeans::emmeans()} + \texttt{pairs()} \\
\S{}7 Ordinal & \texttt{PROC LOGISTIC / link=cumlogit} & \texttt{MASS::polr(method=``logistic'')} \\
\S{}7 Score Test & Proportional Odds Assumption & \texttt{brant::brant()} \\
\S{}7 Binaire & \texttt{PROC LOGISTIC / model Y(event=``1'')} & \texttt{glm(family=binomial)} \\
\S{}7 Stepwise logit & \texttt{PROC LOGISTIC / selection=stepwise} & \texttt{MASS::stepAIC(direction=``both'')} \\
\S{}7 ROC & \texttt{PROC LOGISTIC plots(roc)} & \texttt{pROC::roc()} + \texttt{plot()} \\
\bottomrule
\end{tabular}

\medskip
\noindent\textbf{Packages requis :}
\begin{footnotesize}
\begin{verbatim}
install.packages(c("dplyr", "ggplot2", "car", "MASS", "lmtest", "sandwich",
                   "corrplot", "pROC", "emmeans", "leaps", "vcd", "brant"))
\end{verbatim}
\end{footnotesize}

\bigskip

\begin{footnotesize}
\begin{verbatim}
###########################################################################
#  ANNEXE C -- Code R complet
#  Reproduction fidele des analyses SAS (analyse_etat_sante_esps.sas)
###########################################################################

library(dplyr);   library(ggplot2); library(car);     library(MASS)
library(lmtest);  library(sandwich); library(corrplot)
library(pROC);    library(emmeans); library(leaps);   library(vcd)

# ==================================================================
# 0.  IMPORTATION ET RECODAGE  (equiv. PROC IMPORT + DATA step)
# ==================================================================
esps <- read.csv2("data/base_esps.csv",
                  fileEncoding = "UTF-8", stringsAsFactors = TRUE)
str(esps);  print(table(esps$etat_sante_percu))

esps$score_sante <- dplyr::case_when(
  esps$etat_sante_percu == "Tres mauvais" ~ 1L,
  esps$etat_sante_percu == "Mauvais"      ~ 2L,
  esps$etat_sante_percu == "Assez bon"    ~ 3L,
  esps$etat_sante_percu == "Bon"          ~ 4L,
  esps$etat_sante_percu == "Tres bon"     ~ 5L,
  TRUE ~ NA_integer_
)
esps$etat_sante_percu <- factor(
  esps$etat_sante_percu,
  levels  = c("Tres mauvais","Mauvais","Assez bon","Bon","Tres bon"),
  ordered = TRUE)
esps$sante_mauvaise  <- ifelse(esps$score_sante <= 2, 1L, 0L)
esps$hospit_num      <- ifelse(esps$hospitalisation_12_mois   == "Oui", 1, 0)
esps$handicap_num    <- ifelse(esps$handicap_declare          == "Oui", 1, 0)
esps$renoncement_num <- ifelse(esps$renoncement_soins_12_mois == "Oui", 1, 0)
print(table(as.character(esps$etat_sante_percu), esps$score_sante))

# ==================================================================
# 1.  ANALYSE BIVARIEE EXPLORATOIRE
# ==================================================================
# NOTE : exploratoire uniquement -- selection = cadre theorique.

# 1a. Chi2 + V de Cramer  (equiv. PROC FREQ / chisq measures)
vars_quali <- c("sexe","profession","niveau_education","situation_familiale",
  "couverture_complementaire","hospitalisation_12_mois","tabagisme",
  "consommation_alcool","activite_physique","handicap_declare",
  "renoncement_soins_12_mois","maladies_declarees")
for (v in vars_quali) {
  tbl <- table(esps$etat_sante_percu, esps[[v]])
  cat("\n---", v, "---\n"); print(chisq.test(tbl))
  cat("V de Cramer :", round(vcd::assocstats(tbl)$cramer, 4), "\n")
}

# 1b. ANOVA exploratoire  (equiv. PROC GLM / tukey) -- role indicatif seul
vars_quant <- c("age","revenus_mensuels_euros","imc",
  "depenses_sante_annuelles_euros",
  "nombre_consultations_annuelles","nombre_medicaments_reguliers")
for (v in vars_quant) {
  fit_av <- aov(as.formula(paste(v, "~ etat_sante_percu")), data = esps)
  cat("\n--- ANOVA :", v, "---\n"); print(summary(fit_av))
  print(TukeyHSD(fit_av))
}

# 1c. Kruskal-Wallis  (equiv. PROC NPAR1WAY wilcoxon)
for (v in vars_quant) {
  kt <- kruskal.test(as.formula(paste(v, "~ etat_sante_percu")), data = esps)
  cat(sprintf("%-40s H=%.3f  p=%.4f\n", v, kt$statistic, kt$p.value))
}

# 1d. Spearman -- TEST PRINCIPAL quantitatif  (equiv. PROC CORR spearman)
vars_with <- c("age","revenus_mensuels_euros","imc",
  "depenses_sante_annuelles_euros","nombre_consultations_annuelles",
  "nombre_medicaments_reguliers","nombre_enfants")
for (v in vars_with) {
  r <- cor.test(esps$score_sante, esps[[v]], method = "spearman")
  cat(sprintf("%-40s rho=%.4f  p=%.4f\n", v, r$estimate, r$p.value))
}

# 1e. Boxplots  (equiv. PROC SGPLOT vbox)
for (v in c("age","imc","depenses_sante_annuelles_euros",
            "nombre_consultations_annuelles")) {
  print(
    ggplot(esps, aes(x = etat_sante_percu, y = .data[[v]],
                     fill = etat_sante_percu)) +
      geom_boxplot(alpha = 0.7) +
      scale_fill_brewer(palette = "RdYlGn") +
      labs(x = "Etat de sante percu", y = v,
           title = paste("Distribution de", v)) +
      theme_minimal() + theme(legend.position = "none")
  )
}

# ==================================================================
# 2.  ROBUSTESSE -- REGRESSION LINEAIRE  (equiv. PROC REG)
# ==================================================================
lm_complet <- lm(score_sante ~ age + revenus_mensuels_euros +
  depenses_sante_annuelles_euros + nombre_consultations_annuelles +
  nombre_medicaments_reguliers + imc + nombre_enfants, data = esps)
print(summary(lm_complet))

# Modele reduit -- base THEORIQUE (Grossman, 1972)
lm_reduit <- lm(score_sante ~ age + nombre_consultations_annuelles +
                  nombre_medicaments_reguliers, data = esps)
print(summary(lm_reduit))

# ==================================================================
# 3.  COLINEARITE  (equiv. PROC REG / vif tol collin ; PROC CORR)
# ==================================================================
vif_vals <- car::vif(lm_complet)
print(round(vif_vals, 4)); print(round(1/vif_vals, 4))  # Tolerance
vars_pred <- c("age","revenus_mensuels_euros","depenses_sante_annuelles_euros",
  "nombre_consultations_annuelles","nombre_medicaments_reguliers",
  "imc","nombre_enfants")
cor_mat <- cor(esps[, vars_pred], use = "complete.obs")
print(round(cor_mat, 3))
corrplot::corrplot(cor_mat, method = "color", type = "upper",
  addCoef.col = "black", number.cex = 0.6, tl.cex = 0.75)
sv <- svd(model.matrix(lm_complet))$d
cat(sprintf("Condition number : %.2f  (seuil severe : 30)\n", max(sv)/min(sv)))

# ==================================================================
# 4.  SELECTION AUTO -- VERIFICATION DE COHERENCE SEULEMENT
# ==================================================================
lm_null_sel <- lm(score_sante ~ 1, data = esps)
lm_full_sel <- lm(score_sante ~ age + revenus_mensuels_euros +
  depenses_sante_annuelles_euros + nombre_consultations_annuelles +
  nombre_medicaments_reguliers + imc + nombre_enfants, data = esps)

# 4a. Tous sous-modeles  (equiv. PROC REG / selection=ADJRSQ best=5)
reg_sub <- leaps::regsubsets(
  score_sante ~ age + revenus_mensuels_euros + depenses_sante_annuelles_euros
    + nombre_consultations_annuelles + nombre_medicaments_reguliers
    + imc + nombre_enfants, data = esps, nvmax = 7)
reg_sum <- summary(reg_sub)
cat("R2 ajuste :\n"); print(round(reg_sum$adjr2, 4))
print(reg_sum$which[order(reg_sum$adjr2, decreasing = TRUE)[1:5], ])

# 4b. Forward  (equiv. PROC REG / selection=forward slentry=0.05)
lm_forward <- step(lm_null_sel, direction = "forward",
  scope = list(lower = lm_null_sel, upper = lm_full_sel), trace = 1)
print(summary(lm_forward))

# 4c. Backward  (equiv. PROC REG / selection=backward slstay=0.05)
lm_backward <- step(lm_full_sel, direction = "backward", trace = 1)
print(summary(lm_backward))

# 4d. Stepwise -- VERIFICATION DE COHERENCE SEULEMENT
lm_stepwise <- step(lm_null_sel, direction = "both",
  scope = list(lower = lm_null_sel, upper = lm_full_sel), trace = 1)
print(summary(lm_stepwise))

# ==================================================================
# 5.  VALIDATION DU MODELE LINEAIRE  (robustesse)
# ==================================================================
resid_df <- data.frame(
  resid    = residuals(lm_reduit),  yhat     = fitted(lm_reduit),
  rstud    = rstudent(lm_reduit),   cooksd   = cooks.distance(lm_reduit),
  leverage = hatvalues(lm_reduit))

# 5b. Normalite  (equiv. PROC UNIVARIATE normal qqplot)
qqnorm(resid_df$resid); qqline(resid_df$resid, col = "red")
print(shapiro.test(resid_df$resid))

# 5c. Histogramme
hist(resid_df$resid, breaks = 25, probability = TRUE, col = "steelblue")
curve(dnorm(x, 0, sd(resid_df$resid)), add = TRUE, col = "red", lwd = 2)

# 5d-5e. RStudent x Leverage + Cook
plot(resid_df$leverage, resid_df$rstud,
  xlab = "Leverage", ylab = "RStudent", pch = 19, cex = 0.7)
abline(h = c(-2, 2), lty = 2, col = "red")
threshold_cook <- 4 / nrow(esps)   # 4/500 = 0.008
plot(resid_df$cooksd, type = "h", ylab = "Cook's D")
abline(h = threshold_cook, col = "red", lty = 2)

# 5f. Durbin-Watson  (equiv. PROC REG / dwprob)
print(lmtest::dwtest(lm_reduit))

# 5g. Test de White  (equiv. PROC REG / spec)
print(lmtest::bptest(lm_reduit, studentize = FALSE))

# 5h. Residus partiels  (equiv. PROC REG plots(partialplot))
car::avPlots(lm_reduit)

# 5i. LOESS  (equiv. PROC SGPLOT loess)
for (v in c("age","nombre_consultations_annuelles",
            "nombre_medicaments_reguliers")) {
  print(
    ggplot(esps, aes(x = .data[[v]], y = score_sante)) +
      geom_point(alpha = 0.3, size = 1) +
      geom_smooth(method = "loess", se = TRUE, color = "red") +
      labs(x = v, y = "score_sante") + theme_minimal()
  )
}

# 5j. Observations aberrantes
idx_out <- which(abs(resid_df$rstud) > 2 | resid_df$cooksd > threshold_cook)
print(cbind(esps[idx_out, c("id","age","score_sante")],
            round(resid_df[idx_out, ], 4)))

# ------------------------------------------------------------------
# 5k. K-FOLD (k=10) -- modele ordinal  (equiv. macro %kfold_ordinal)
# ------------------------------------------------------------------
set.seed(42);  k <- 10L;  n <- nrow(esps)
folds <- rep(NA_integer_, n)
for (lev in levels(esps$etat_sante_percu)) {
  idx   <- which(as.character(esps$etat_sante_percu) == lev)
  idx_s <- sample(idx)
  folds[idx_s] <- (seq_along(idx_s) - 1L) %% k + 1L
}
concordance_folds <- numeric(k)
for (i in seq_len(k)) {
  fit_fold <- MASS::polr(
    etat_sante_percu ~ age + nombre_consultations_annuelles +
      nombre_medicaments_reguliers + sexe + hospitalisation_12_mois +
      activite_physique + handicap_declare,
    data = esps[folds != i, ], Hess = TRUE, method = "logistic")
  pred <- predict(fit_fold, newdata = esps[folds == i, ])
  concordance_folds[i] <- mean(pred == esps$etat_sante_percu[folds == i])
  cat(sprintf("  Fold %2d : %.4f\n", i, concordance_folds[i]))
}
cat(sprintf("Concordance moy : %.4f  sd : %.4f\n",
            mean(concordance_folds), sd(concordance_folds)))

# ==================================================================
# 6.  ROBUSTESSE -- ANCOVA  (equiv. PROC GLM/ss3 + LSMEANS/tukey)
# ==================================================================
esps$sexe <- relevel(esps$sexe, ref = "Homme")
esps$hospitalisation_12_mois <- relevel(esps$hospitalisation_12_mois, ref="Non")
esps$tabagisme <- relevel(esps$tabagisme, ref = "Non-fumeur")
esps$activite_physique <- relevel(esps$activite_physique, ref = "Sedentaire")
esps$handicap_declare <- relevel(esps$handicap_declare, ref = "Non")
esps$couverture_complementaire <- relevel(esps$couverture_complementaire,
                                           ref = "Mutuelle privee")

# 6a. ANCOVA complete
lm_ancova_complet <- lm(score_sante ~ age + revenus_mensuels_euros + imc +
  depenses_sante_annuelles_euros + nombre_consultations_annuelles +
  nombre_medicaments_reguliers + sexe + profession + couverture_complementaire
  + hospitalisation_12_mois + tabagisme + activite_physique +
  handicap_declare + maladies_declarees, data = esps)
print(car::Anova(lm_ancova_complet, type = "III"))

# 6b. ANCOVA reduite
lm_ancova_reduit <- lm(score_sante ~ age + nombre_consultations_annuelles +
  nombre_medicaments_reguliers + sexe + hospitalisation_12_mois +
  activite_physique + handicap_declare, data = esps)
print(car::Anova(lm_ancova_reduit, type = "III"))
print(summary(lm_ancova_reduit))
for (fac in c("sexe","hospitalisation_12_mois","activite_physique",
              "handicap_declare")) {
  cat("\n---", fac, "---\n")
  print(pairs(emmeans::emmeans(lm_ancova_reduit, fac), adjust = "tukey"))
}

# 6c. Interaction Age x Sexe
lm_ia <- lm(score_sante ~ age + sexe + age:sexe, data = esps)
print(car::Anova(lm_ia, type = "III"))

# ==================================================================
# 7.  MODELE PRINCIPAL -- REGRESSION LOGISTIQUE ORDINALE
# ==================================================================
# MASS::polr() = equiv. exact PROC LOGISTIC / link=cumlogit

# 7a. Modele complet
polr_complet <- MASS::polr(
  etat_sante_percu ~ age + revenus_mensuels_euros + imc +
    nombre_consultations_annuelles + nombre_medicaments_reguliers +
    depenses_sante_annuelles_euros + sexe + hospitalisation_12_mois +
    tabagisme + activite_physique + handicap_declare + couverture_complementaire,
  data = esps, Hess = TRUE, method = "logistic")
print(summary(polr_complet))
cat("OR :\n"); print(round(exp(coef(polr_complet)), 4))
print(round(exp(confint(polr_complet)), 4))

# Pseudo-R2 de Nagelkerke  (equiv. option rsquare de SAS)
ll_null_p  <- logLik(MASS::polr(etat_sante_percu ~ 1, data = esps,
                                 method = "logistic", Hess = FALSE))
ll_model_p <- logLik(polr_complet);  n_obs <- nrow(esps)
r2_nag <- (1 - exp((2/n_obs)*(as.numeric(ll_null_p)-as.numeric(ll_model_p)))) /
          (1 - exp((2/n_obs)*as.numeric(ll_null_p)))
cat(sprintf("Pseudo-R2 Nagelkerke : %.4f  AIC : %.2f  BIC : %.2f\n",
            r2_nag, AIC(polr_complet), BIC(polr_complet)))

# 7b. Test de proportionnalite des odds (equiv. Score Test de SAS)
if (requireNamespace("brant", quietly = TRUE)) {
  print(brant::brant(polr_complet))
} else cat("install.packages('brant') pour le test de proportionnalite\n")

# 7c. Modele reduit -- MODELE PRINCIPAL FINAL  (base theorique)
polr_reduit <- MASS::polr(
  etat_sante_percu ~ age + nombre_consultations_annuelles +
    nombre_medicaments_reguliers + sexe + hospitalisation_12_mois +
    activite_physique + handicap_declare,
  data = esps, Hess = TRUE, method = "logistic")
print(summary(polr_reduit))
cat("OR :\n"); print(round(exp(coef(polr_reduit)), 4))
print(round(exp(confint(polr_reduit)), 4))
cat(sprintf("AIC=%.2f  BIC=%.2f\n", AIC(polr_reduit), BIC(polr_reduit)))

# 7d. Logistique binaire (robustesse)
logit_bin <- glm(sante_mauvaise ~ age + revenus_mensuels_euros + imc +
  nombre_consultations_annuelles + nombre_medicaments_reguliers +
  sexe + hospitalisation_12_mois + tabagisme + activite_physique +
  handicap_declare, data = esps, family = binomial)
print(summary(logit_bin))
cat("OR :\n"); print(round(exp(coef(logit_bin)), 4))
pred_bin <- ifelse(predict(logit_bin, type = "response") >= 0.5, 1, 0)
print(table(Observe = esps$sante_mauvaise, Predit = pred_bin))
cat(sprintf("Taux : %.1f%%\n", mean(pred_bin == esps$sante_mauvaise)*100))

# 7e. Coherence stepwise -- VERIFICATION SEULEMENT
polr_step <- MASS::stepAIC(polr_complet, direction = "both", trace = 1)
print(summary(polr_step))

# 7f. Courbe ROC  (equiv. PROC LOGISTIC plots(roc))
logit_bin_red <- glm(sante_mauvaise ~ age + nombre_consultations_annuelles +
  nombre_medicaments_reguliers + sexe + hospitalisation_12_mois +
  activite_physique + handicap_declare, data = esps, family = binomial)
roc_obj <- pROC::roc(esps$sante_mauvaise,
  predict(logit_bin_red, type = "response"), quiet = TRUE)
cat(sprintf("AUC = %.4f\n", pROC::auc(roc_obj)))
plot(roc_obj, main = paste("Courbe ROC -- AUC =", round(pROC::auc(roc_obj),4)),
     col = "steelblue", lwd = 2)
abline(a = 0, b = 1, lty = 2, col = "grey")

# 7g. Matrice de confusion  (equiv. ctable de SAS)
pred_r <- ifelse(predict(logit_bin_red, type = "response") >= 0.5, 1, 0)
cf <- table(Observe = esps$sante_mauvaise, Predit = pred_r)
print(cf)
cat(sprintf("Sensibilite : %.1f%%  Specificite : %.1f%%  Global : %.1f%%\n",
  cf[2,2]/(cf[2,2]+cf[2,1])*100,
  cf[1,1]/(cf[1,1]+cf[1,2])*100,
  (cf[2,2]+cf[1,1])/sum(cf)*100))

###########################################################################
# FIN DE L'ANNEXE C
###########################################################################
\end{verbatim}
\end{footnotesize}

% --- Bibliographie ------------------------------------------------------------
\printbibliography[heading=bibintoc,title={Bibliographie}]

\end{document}
